%! TeX program = lualatex
\documentclass[12pt]{report}

% Language and Encoding
\usepackage{fontspec}
\usepackage[bidi=basic]{babel}



\defaultfontfeatures{ Scale=MatchUppercase,
                      Ligatures=TeX,
                      Renderer=HarfBuzz }

% German
\babelprovide[import, main]{german}

% Greek
\babelprovide[import, onchar=ids fonts]{greek}
\babelfont[greek]
           {rm}
           [Path = ./static/fonts/GalatiaSIL-2.1-web/,
           Extension = .ttf,
           Renderer=HarfBuzz,
           Language=Default ]
          {GalSILR}

% Hebrew
\babelprovide[import=he, onchar=ids fonts]{hebrew}
\babelfont[hebrew]
           {rm}
           [Path = ./static/fonts/EzraSIL2.51/,
           Extension = .ttf,
           Renderer=HarfBuzz,
           Language=Default ]
          {SILEOTSR}


% graphics
\usepackage{graphicx}
\graphicspath{ {./static/images} }

% Bibtex
\usepackage{csquotes}
\usepackage[backend=biber,style=alphabetic,citestyle=alphabetic]{biblatex}
\addbibresource{./static/Theologie_Bibliography.bib}
% No S. in front of postnote, because I can't cite non-standard books otherwise
\DeclareFieldFormat{postnote}{#1}

% urls
\usepackage{url}
\setcounter{biburllcpenalty}{2000}
\setcounter{biburlucpenalty}{2000}

% Citing the Bible
\input{./static/bible-citation.tex}

% Example Numbering
\usepackage{ifthen}
\usepackage{environ}
\usepackage{letltxmacro}
\usepackage{amssymb}

% \usepackage{gb4e}

% root-sign
\newcommand{\lroot}[1]{\raisebox{1mm}{$\sqrt{}$}#1}

\newcommand{\MT}{$\mathfrak{M}$ }
\newcommand{\LXX}{$\mathfrak{G}$ }
\newfontfamily\DoulosSIL{Doulos SIL}
\newcommand{\SamPent}{{\DoulosSIL ^^^^214f}}


% MISC
\usepackage{hyperref}
\hypersetup{
    colorlinks=false,
    pdfborder={0 0 0},
    citecolor=red,
    linkcolor=blue
}
\usepackage{dsfont}
\usepackage{enumerate}
 
\usepackage{imakeidx}

% Special Theorem style
\usepackage{amsthm}
\newtheoremstyle{aspektd}% name of the style to be used
  {\topsep}% measure of space to leave above the theorem. E.g.: 3pt
  {\topsep}% measure of space to leave below the theorem. E.g.: 3pt
  {\mdseries}% name of font to use in the body of the theorem
  {0pt}% measure of space to indent
  {\bfseries}% name of head font
  { ---}% punctuation between head and body
  { }% space after theorem head; " " = normal interword space
  {\thmname{#1}\thmnumber{ #2}\textnormal{\thmnote{ (#3)}}}

\theoremstyle{aspektd}
\newtheorem{aspekt}{Aspekt}

% special words
\newcommand{\kipper}{כִּפֶּר (kipper)}

\title{UNPUBLISHED - Penal Substitutionary Atonement - UNPUBLISHED}
\date{}
\author{Jonathan Schleucher}

\makeindex

\begin{document}

{\let\newpage\relax\maketitle}

\abstract{
TODO
}

\tableofcontents

\section*{Vorwort}
Während ich nach und nach die einzelnen Kapitel für diesen Artikel schreibe steht in jedem unfertigen Abschnitt eine minimale Zusammenfassung
meiner Sicht, ohne oder mit nur minimaler Begründung.
Penal Substitutionary Atonement werde ich durchweg mit PSA abkürzen.

tl;dr des Vorworts:
Hierfür werde ich argumentieren:
\begin{itemize}
\item Gott kann unabhängig von Jesu Tod Sünde vergeben (und tut das auch)
\item Jesus wurde am Kreuz nicht vom Vater an unserer Stelle für unsere Sünden bestraft
\item Jesus ist als Opfer für unsere Sünden gestorben; das bedeutet aber etwas anderes als PSA\footnote{
	Siehe Kapitel \ref{chap-hebrews} für eine ausführliche Erklärung, wie dieses Opfer im neuen Testament explizit definiert wird}
\end{itemize}

\pagebreak
\part{Das levitische System}
Um ein gutes Verständnis von Jesus wirken als Opfer zu erlangen müssen wir zunächst einige Grundlagen im Alten Testament legen.
In diesem Teil werde ich für folgende Aussagen argumentieren:
\begin{enumerate}[(1)]
\item Kein Opfer im gesamten levitischen System kann als `substitution' des Opertieres an Stelle des Opfernden verstanden werden.
\item Bei keinem Opfer im levitischen System ist der `Tod' des Opfertieres zentral. Ganz im Gegenteil wird das Schlachten des Opfertieres rituell zu einem `nicht-töten' rekonzeptualisiert.
\item Es gibt kein Ritual oder Opfer im levitischen System, dass `Sünden vergeben' hat. Insbesondere sind die üblichen deutschen Übersetzungen `Sündenopfer', und `entsühnung' grundlegende Missverständnisse der Ziele und Limitierungen des levitischen Systems.
\end{enumerate}

\chapter{Grundlegende Konzepte des levitischen Systems}
Bevor wir die theologische Signifikanz des Opfersystems verstehen können, ist es nötig einige fundamentale Grundlagen des levitischen Systems zu verstehen.
Während wir dieses Grundverständnis erlangen ist es auch nötig, einige sehr übliche Wörter in den deutschen Übersetzungen des Alten Testaments aus unserem Sprachgebrauch zu entfernen.

\section{Heilig und Gemein}
Eine der wichtigsten Distinktionen, die das levitische System untermauert ist die Unterscheidung zwischen קֹדֶשׁ (qodeš, heilig) und חֹל (ḥol, gemein).
Heiliges ist (zu verschiednene Graden) abgesondert für einen speziellen Nutzen in Gottes Nähe.
Zunächst wertungsfrei sind gemeine Dinge oder Menschen solche, die nicht heilig sind.

Beispiele für heilige Dinge sind: das innere des Zeltes der Begegnung und des Tempels,
Gegenstände im Heiligtum,
die Kleider der Priester\bibfoot{Exodus}{40}{13},
die Priester selbst\bibfoot{Exodus}{40}{13},
das Fleisch verschiedener Opfer
\footnote{\bibverse{Leviticus}{10}{17}. Beachte, dass dieses Opferfleisch sogar hochheilig ist. Heiligkeit im levitischen System hat verschiedene Stufen.}
.

\section{Rein und Unrein}
Die zweite, und von Heiligkeit/Gemeinheit \textit{völlig} unabhängige Dimension ist die Unterscheidung zwischen טָחוֺר (ṭaḥor, rein) und טָמֵא (ṭame, unrein).
Im levitischen System gibt es verschiedene Arten von Reinheit, die es zu unterscheiden gilt.
Alle diese verschiednen Arten verwenden dieselben hebräischen Worte, sind aber in der Praxis leicht voneinander zu unterscheiden.
\begin{enumerate}[(a)]
	\item Jedes Ding und jeder Mensch ist zu einem bestimmten Zeitpunkt immer entweder \textit{rituell} rein oder \textit{rituell} unrein.
	Rituelle Unreinheit ist durch Berührung übertragbar\bibfoot{Numbers}{19}{22}.
	\item Manche Handlungen sind verboten, und dieses Verbot wird mit Reinheitssprache ausgedrückt.
	\footnote{In diese Kategorie fallen insbesondere die unreinen Tiere. Diese Tiere sind nicht grundsätzlich rituell unrein, denn Berührung eines unreinen Tieres überträgt keine rituelle Unreinheit. Ihr verzehr ist absolut verboten, und dieses Verbot drückt beispielsweise \bibverse{Leviticus}{11}{4} als unreinheit aus.}
	\item Moralisch verwerfliche Taten werden als moralische Unreinheit bezeichnet.
		\footnote{Wiederum ist offensichtlich, dass diese Art von Unreinheit unabhängig von ritueller Unreinheit ist. Das berühren eines rituell unreinen Menschen überträgt Unreinheit, das berühren eines moralisch unreinen Menschen (eines Sünders) überträgt keine rituelle (oder gar moralische) Unreinheit.}
\end{enumerate}
Es ist essentiell zu bemerken, dass heilig/unheilig und rein/unrein zwei verschiedene Dimensionen sind.
Ein Priester ist immer heilig (abgesondert für einen bestimmten Dienst für Gott), kann aber beispielsweise durch die Berührung einer Leiche unrein werden.
Das Beispiel aus \bibverse{Leviticus}{21}{1-3} ist instruktiv:
(a) Ein Priester kann sich Leichen-Unreinheit zuziehen wie jeder andere Mensch (b) er darf das im Normalfall nicht tun und (c) wenn es um Blutsverwandte geht ist dieses Verbot aufgehoben.
Wir sehen hieran die folgenden essentiellen Punkte, die wir über zwingend über rituelle Reinheit wissen müssen, bevor wir das Neue Testament lesen:
\begin{enumerate}[(1)]
	\item Sich rituelle Unreinheit zuzuziehen ist \textit{keine} Sünde.
	\item In manchen speziellen Situationen ist das zuziehen ritueller Unreinheit Sünde, und zwar \textit{aus göttlichem Befehl}. Diese Situation ist nicht der Normalzustand.
\end{enumerate}
Dass diese Punkte wahr sind sehen wir genauso auch an \bibverse{Leviticus}{15}{16-18}.
Es gibt keine Kinder ohne rituelle Unreinheit\footnote{Ganz zu schweigen von der Unreinheit der Frau nach der Geburt [\bibverse{Leviticus}{12}{}], oder der Unreinheit der Frau durch ihre Menstruation [\bibverse{Leviticus}{15}{19ff}.},
und es ist absolut normal, dass Menschen häufig rituell unrein werden.

Rituelle Unreinheit wird nur dann problematisch, wenn außerdem heilige Objekte im Spiel sind.
Eines der primären Probleme, dass das levitische System löst ist, dass heilige Objekte niemals mit Unreinheit in berührung kommen dürfen.
Wenn also eine Frau in den ersten sieben Tagen nach dem Beginn ihrer Menstruation (und dementsprechend unrein) ist, so darf sie nicht mit heiligen Gegenständen in Kontakt kommen, und nicht vom Fleisch von Dankopfern essen.
Es gibt keine weiteren Restriktionen oder Einschränkungen, die gewöhnliche rituelle Unreinheit nach sich zieht.
Eine Person die nicht vorhatte am Tempelkult in Jerusalem teilzunehmen hat also kaum einen Grund, viel über rituelle Unreinheit nachzudenken.
Ein rituelles Bad am Tag reicht aus, um den überwältigenden Großteil der rituellen Reinheitsgebote zu erfüllen.

Andere rituelle Unreinheiten können nicht durch ein einfaches Bad entfernt werden. Dazu zählt beispielsweise das berühren von Leichen\bibfoot{Numbers}{19}{}.
Die rituelle Unreinheit wird auch in solchen Fällen durch nicht-Opfer und Zeit entfernt, namentlich rituelle Waschungen, oder andere Riten (wie das besprengen der Person mit einer bestimmten Flüssigkeit).

Es gibt besonders große Fälle von ritueller Unreinheit, nämlich
Unreinheit durch (a) Aussatz\bibfoot{Leviticus}{14}{} (b) die Geburt eines Kindes\bibfoot{Leviticus}{12}{}, und (c) Blutfluss von Mann oder Frau\bibfoot{Leviticus}{15}{}.
Der entscheidende Unterschied zwischen großer und gewöhnlicher ritualler Unreinheit ist, dass große rituelle Unreinheit sich auf das Heiligtum überträgt, selbst ohne örtliche Proximität.
\Footcite[78f]{lamb-of-the-free}

Das bedeutet zum Beispiel, dass rituelle Unreinheit durch die Geburt eines Kindes auf das Heiligtum übertragen wird - das heilige kommt mit unreinem in Berührung, und für dieses Problem braucht es eine Lösung.
Wir kommen jetzt zu einem der wichtigsten Punkte, die wir über die Opfer im levitischen System verstehen müssen: Was ist `Entsühnung' und was bewirkt ein `Sündopfer'?

\section{\label{kipper}Dekontamination}
Die Lösung im levitischen System für Unreinheit, die auf das Heiligtum übertragen wurde ist die `Entsühnung', die ein `Sündopfer' bringen kann.

Das erste dieser Wörter übersetzt häufig das hebräische \kipper.
Die deutschen Glossen in einigen gängigen Übersetzungen sind:
`Sühnung erwirken'\footnote{\Cite{elb}, \Cite{meng}, \Cite{neü}, \Cite{ngü}, \Cite{slt}}, `Sühnung vollziehen'\Footcite{lut}, `Wiedergutmachung schaffen'\Footcite{nlb}, und `Versöhnung erwirken'\Footcite{eü}.

Alle diese Übersetzungen sind völlig falsch.\footnote{
Manchmal ist Klarheit nötig. Das Referenzwerk für die gesamte Diskussion um das levitische System ist \Cite{milgrom-leviticus-1}, siehe dort 255f für die Diskussion über \kipper.}

Das hebräische \kipper
\footnote{\lroot{כפר} im Piel}
hat als primäre Bedeutung `überdecken', wird aber insbesondere im levitischen System restriktiver verwendet: als `dekontaminieren' oder `reinigen' heiliger Gegenstände.\Footcite[4]{lamb-of-the-free}
Als anschauliche Begründung für diese Tatsache kann ein Vergleich des
`Sündopfers' für unwissentlich begangene Sünden des Volkes, \bibverse{Leviticus}{4}{20},
und des finale `Sündopfers' nach einer Geburt, \bibverse{Leviticus}{12}{7}, helfen.
Im ersten Fall lesen wir ``und der priester soll für sie kipper-wirken, und ihnen wird vergeben werden'', im zweiten ``und er soll kipper-wirken für sie, und sie wird rein sein''.
Wir sehen also, dass kipper sowohl zeitgleich mit `vergeben' als auch mit `rein werden' stehen kann.
Im ersten Fall gab es eine Sünde (und die damit verbundene verunreinigung des Heiligtums), im zweiten Fall gab es eine große rituelle Unreinheit (die sich auf das Heiligtum übertragen hat).
Im ersten Fall geschieht kipper und Vergebung, im zweiten Fall kipper und rein-Werden.
Das bedeutet, dass kipper eine Bedeutung haben muss, die völlig ohne Relation zu Sünde formulierbar ist (alternativ wäre das Gebähren eines Kindes eine Sünde).
Die Bedeutung die übrig bleibt\footnote{Also die maximale invariante (lexische) Semantik von \kipper, die diese beiden Stellen ohne einfügen eines homonym auflösen kann} ist ``Verunreinigung entfernen''.
Aus diesem Grund werde ich kipper durchgängig mit Dekontamination übersetzen und ab hier nur noch von `Entsühnung' sprechen, wenn das entfernen von Sünde tatsächlich ist was ich meine.

Eine weiteres fundamentales Missverständnis ist, dass das Blut von Opfern auf Menschen aufgetragen wird, um sie von Unreinheit und Schuld zu reinigen.
Im nächsten Kapitel werden wir alle Blutrituale ins Auge nehmen, bei denen Blut auf Menschen aufgetragen wird und werden sehen, dass kein Blut das als Teil eines Rituals Menschen berührt
\footnote{Menschen berühren soll. Es gibt Fälle, wo Opferblut von dekontaminierenden Opfern aus Versehen Dinge außer die heiligen Objekte berührt, die dekontaminiert werden sollen; das ist nicht beabsichtigt und problematisch.}
von einem dekontaminierenden Opfer stammt.
Aber das Verb kipper im gesamten Levitischen Text hat \textit{nie} eine Person als direktes Objekt.
Es wird `für eine Person dekontaminiert', aber das direkt Objekt dieser Dekontamination ist immer ein Objekt\footnote{oder elliptisch}, und niemals eine Person. Es wird `für eine Person ein Objekt dekontaminiert'.
Dieses Phänomen werden wir genauer untersuchen, wenn wir über das Reinigungsopfer (`Sündopfer') reden.


\chapter{Nicht dekontaminierende Opfer}
Was ist die primäre Aufgabe des levitischen Systems, und speziell des Opfersystems?
Wer in einer einfachen Version von PSA aufgewachsen ist könnte denken, dass die primäre Funktion der Opfer im Alten Testament das entfernen von Sünde und hindeuten auf den viel größeren stellvertretenden Opfertod des Messias ist.
Wir haben bereits im vergangenen Abschnitt festgehalten, dass kipper (das hinter den deutschen Übersetzungen `entsühnen' steht) nichts direkt mit Sünde oder Vergebung zu tun hat.
Was also ist dann die Funktion des Opfersystems?

Im großen ganzen sind das Opfersystem und das levitische System im ganzen eingesetzt, damit Israel mit ihrem Gott (der in ihrer Mitte wohnt) in Gemeinschaft stehen können.
So beginnt und endet jeder Tag mit einem regelmäßigen Brandopfer, das in \bibverse{Numbers}{28}{} eingesetzt wird, als "Gottes Speise", "ihm zum wohlgefälligen Geruch".
Dieses Brandopfer zieht die Aufmerksamkeit Gottes auf sich - ein Opfer, dass vollständig verbrannt wird und nur deshalb existiert, weil Gott es so wollte.
Neben den Brandopfern spielen die Heilsopfer eine große Rolle -
bei jedem der großen Jahresfeste werden Heilsopfer gebracht, um an die vergangenen Wohltaten Gottes zu erinnern, und auch um sein persönliches Dank an Gott auszudrücken hat ein Israelit ein Heilsopfer gebracht.
Die Heilsopfer sind speziell in der Hinsicht, dass ihr Fleisch unter anderem von der Person gegessen wird, die das Opfer bringt; es ist ein gemeinsames Mahl mit Gott:
ein drittel geht in Flammen vor Gott auf, ein drittel gehört den Priestern (die kein Land und deswegen keine eigenen Herden haben), ein drittel isst der Opfernde selbst.
Hier ist also in gröbsten Zügen die Funktion des Opfersystems:
(a) Brandopfer werden gebracht um Gott "anzuziehen", ihm als Wohlgeruch
(b) Heilsopfer werden dargebracht um Gott zu loben und ihm zu danken
(c) durch Unreinheit (moralisch und rituell) wird heiliger Ort verunreinigt; die dekontaminierenden Opfer sorgen dafür, dass diese Verunreinigung vom Heiligtum entfernt wird damit Gottes Heiligkeit weiterhin dort wohnen kann.


\section{Brandopfer und Heilsopfer}
- Beschreibung von Brandopfer und Heilsopfer
- Heilsopfer werden gegessen - nur nicht-dekontaminierende Opfer können gegessen werden
- Fußnote oder kurzer Satz: zum Teil hat olah wohl kipper-wirkung (cf. lev 1).
\section{Die Opfer zu Passah, Bundesschluss, und Bundeserneuerung}
- Brandopfer + Heilsopfer
- Erinnerung an Gottes Erlösung aus Ägypten
- Dankbarkeit für Gottes Bund
- Rededikation und neues Einladen Gottes nach Zeiten großer Sünde
\section{Blutrituale mit menschlichem Objekt}
- Metaphysischer Übergang
- Blut ist IMMER von nicht-dekontaminierenden Opfern

\chapter{Dekontaminierende Opfer}
- Funktion: entfernen von Unreinheit an heiligen Objekten
- Menschen werden IMMER UND AUSSCHLIEßLICH durch Waschung, Zeit oder andere Rituale gereinigt
\section{Das Reinigungsopfer}
- hat nichts mit Sünde zu tun
- Leben ist im Blut
- Tod des Tieres ist dafür da, um an Blut zu kommen
	- im vordergrund steht das Blutritual und was mit dem Fleisch des Tieres passiert
- Opfer im AT ist kein rituelles Schlachten; Schlachten ist eine notwendigkeit, um Blut und Fleisch zu haben
\section{Das Reparationsopfer}
- Unterart des Reinigungsopfers
- muss ich mir nochmal speziell anschauen
\section{Der Tag der Dekontamination}
- Ziel: Dekontamination des Heiligtums von innen nach außen
	- keine substitution; Opfer für das Volk, nicht Substitutionell an ihrer Stelle
- der Sündenbock kommt später
\section{Fallbeispiel: Reinigung nach der Geburt}
- Reinigung der Person durch Zeit und Waschung
- danach (Wenn die Person rein ist) hattat, um das Heiligtum zu kippern und "die Person wird rein"
- Mishna dazu zitieren; im prinzip aus lamb-of-the-free abschreiben

\chapter{Rituale, die keine Opfer sind}
% ich denke das kann man weglassen, weil die relevanten Punkte schon angesprochen sind
%\section{Reinigungsrituale für kleine rituelle Unreinheit}
%\section{Reinigungsrituale für große rituelle Unreinheit}
\section{\label{metaphysische-transformation}Der Bundesschluss und Blutreinigungsrituale}
- Verbindung zwischen Volk und Gott ausgedrückt als besprenklung
	- kein kipper, keine Vergebung
\section{Der Sündenbock}
- Sünde wird nicht kippert, sie wird übertragen und entfernt
- OPFER NEHMEN KEINE SÜNDE WEG
	- an dem einen Punkt, wo sünden aus dem Lager entfernt werden ist das KEIN OPFER
\section{Das erste Passah}
- kein opfer
- das muss ich mir nochmal genauer anschauen - was geht hier vor??
	- iwas mit "welchem Gott folgt ihr nach?", weil das Lamm ein Idol von Atum ist?? steckt da was dahinter?

%% ggf. auch direkt oben eingebaut
\chapter{Was Opfer nicht sind}
\section{Substitution}
das Handauflegungsritual
\section{\label{opfer-als-nicht-tod}Tod, Strafe, und Leiden}
Rituelle Rekonzeption des Todes
Der Tod des Tieres ist ein Mittel Zum zweck: um an Blut und Fleisch zu kommen



\pagebreak
\part{Sünde und Vergebung}
\chapter{Verunreinigung des Landes}
tl;dr:
Manche Sünden sind so groß, dass sie nicht bloß das Heiligtum verunreinigen, sondern dass sie `das Land' selbst verunreinigen.\footnote{
Das sind: Götzendienst, Sexuelle Vergehen, und Blutvergießen.}
Für diese Sünden gibt es keine Reparationsopfer und auch ihre Verunreinigung kann nicht einmal durch Säuberungsopfer oder Yom Kippur entfernt werden, weil
diese Opfer nur das Heiligtum dekontaminieren.
Die einzigen Möglichkeiten, mit diesen Opfern umzugehen ist, die Todesstrafe zu vollziehen.
Wenn das Land zu sehr verunreinigt worden ist, folgt das Exil, damit das Land Ruhe hat und durch diese verstrichene Zeit wieder rein werden kann (in Analogie zur Reinigung
der meisten rituellen Unreinheiten die unter anderem durch Zeit geschieht).

\chapter{\label{chap:bundesfluch}Der Lohn der Sünde und der Fluch des Bundes}
tl;dr:
`Der Lohn der Sünde ist der Tod' und das bedeutet, dass Sünde unweigerlich zum Tod führt.
Tod ist nicht immer eine Strafe Gottes.
Manchmal tötet Gott Menschen als Strafe für ihre Sünden, aber nicht jeder Tod ist Gottes Strafe; Tod ist die natürliche Folge der Sünde.
Tod ist `die Macht des Teufels', und `der letzte Feind, der vernichtet wird'.
Die Sünde verwendet das perfekte Gesetz, das zum Leben führen sollte, um stattdessen Tod hervorzurufen.


% TODO Frage:
% Ist Die "Strafe des Gesetzes" ausschlißlich eine natürliche Konsequenz, genau so wie Tod eine natürliche Konsequenz der Sünde ist?
% - insb. interessant für Gal 3:13
% Die Idee wäre dann:
% "Das Gesetzt bringt leben dem der es befolgt. Durch die Sünde sind Menschen nicht in der Lage es zu halten. Als Konsequenz für das Nicht-Halten kommt Fluch. Jesus besiegt die Sünde, indem er die Konsequenz des Gesetztes trägt, uns reinwäscht und seine immortalität anzieht, um uns von der Sünde zu heilen."
% 
\chapter{\label{chap-fiat}Vergebung durch Fiat}
tl;dr:
Gott vergibt einfach so - meist nachdem die Person Buße getan hat.
Gott muss Sünde nicht bestrafen.
Jesus vergibt in seinem Leben einfach so.
Nie hat sich jemand vor Jesu Tod gefragt wie es möglich ist, dass Gott Sünden vergeben kann ohne sie zu bestrafen.
Gott entscheidet sich durch Fiat, Sünden zu vergeben.

\chapter{\label{chap:prophetische-hoffnung}Die prophetische Hoffnung}
tl;dr:
Die Propheten leben kurz vor oder während des Exils, also in der Zeit in der Israel unter dem Fluch des Gesetzes steht.
Ihre Hoffnung ist nicht, dass es ein endgültiges, finales Säuberungsopfer geben wird, mit dem Gott alle Sünden abwaschen kann.
Stattdessen ist die Prophetische Hoffnung konsistent so beschrieben:
(a) Gott wird Israel aus Gnade wieder annehmen und ihnen vergeben
(b) Gott wird die Verunreinigung der Sünden durch eine Waschung im Heiligen Geist von den Menschen entfernen
(c) Gott wird die Menschen grundsätzlich verändern, um sie zu befähigen, Gott zu dienen.
Weiterhin weiten die Propheten die Kategorie der Sünden, die das Land verunreinigen auf, indem sie Ökonomische Unterdrückung als eine Form von Blutvergießen beschreiben.

% Israel noch (teilweise / spirituell) im Exil (2. Tempel Texte zitieren)
% 4Q504 5:15-16; Isa 44:1-3 :: Relation zwischen dem Ende des Exils und der Ausgießung des heiligen Geistes

\chapter{Der Messias und die Kräfte des Todes}
tl;dr:\footnote{
Dieses Kapitel ist eine kurze Zusammenfassung von \Cite{jesus-and-the-forces-of-death}}
Jesus tritt in den Evangelien auf und kommt immer wieder in Kontakt mit den verschiedenen Formen ritueller Unreinheit.
Anstatt die Menschen nur rituell zu reinigen, heilt er stattdessen die Quelle der Unreinheit - den Aussatz, den Blutfluß, den Tod an sich.
Alle Quellen der rituellen Unreinheit wurden als Kräfte des Todes aufgefasst, und überall wo der Messias mit den Kräften des Todes in Kontakt kommt besiegt er
sie durch sein überfließendes Leben und seine Heiligkeit.
Das große Finale seines Wirkens erwirkt Jesus durch seinen Tod und seine Auferstehung.
Hier kommt er mit dem Tod an sich in Kontakt, und wieder ist sein göttliches Leben dem Tod überlegen - Jesus vernichtet den Tod durch sein Leben.

\chapter{Der Tod des Messias und der neue Bund}
tl;dr:
Jesus selbst setzt das Abendmahl als eines der zentralen Rituale für seine Nachfolger ein.
Dieses Mahl ist explizit Jesus als neues Passahlamm, und die Erinnerung an einen neuen Bundesschluss.
Jesus vergießt sein Blut um diesen Bund einzuleiten und seine Nachfolger dadurch zu reinigen und zu heiligen.
Jesu Wirken erfüllt genau die prophetische Hoffnung, und beendet das Exil endgültig.

\chapter{\label{chap-hebrews}Der Messias und sein Opfer im Hebräerbrief}
Nachdem die Evangelien Jesu Tod als Einläuten des neuen Bundes, die Erfüllung der prophetischen Hoffnung, und seinen Sieg über die Kräfte des Todes dargestellt haben wenden wir uns jetzt dem Hebräerbrief zu.
Hier haben wir den längsten Abschnitt des neuen Testaments, der sich in detaillierter Art und Weise mit der Bedeutung von Jesu Opfer in Bezug auf die Opfer im levitischen System auseinandersetzt.
Es ist als allererstes sehr klar, dass Jesus sich selbst als Opfer dem Vater darbringt, und dass dieses Opfer dafür sorgt, dass Menschen Vergebung
für ihre Sünden erhalten. Dieser Punkt steht nicht zur Debatte.
Wenn wir den Hebräerbrief und seine Interaktion mit Jesus als Opfer verstehen, dann wird uns das ein klares Bild davon geben, wie PSA nicht nötig ist
um mit vollem Herzen sagen zu können `Jesus ist das Opferlamm, das zur Vergebung unserer Sünden geschlachtet wurde.'

Im gesamten Hebräerbrief legt der Author dar, wie Jesus das Alte Testament erfüllt, bloß in größerer, besserer, vollständigerer Form.
Da wir uns speziell mit dem levitischen Opfersystem und Jesus als Opfer auseinandersetzen beginnt die interessante Diskussion für unsere Frage in Kapitel 7.
Hier begründet der Author (a) warum Jesus ein Priester sein kann, obwohl er nicht von Levi abstammt (b) warum er sogar besser ist als die irdischen Priester.
Beides wird begründet, da Jesus ein Hoher Priester in der Ordnung Melchisedeks ist - dem prototypischen Priester Gottes außerhalb des jüdischen Volkes.
Kapitel 8 führt das Argument weiter: Wenn Jesus ein Hoher Priester ist, der besser ist als die irdischen Priester, dann setzt er einen Dienst und einen Bund ein, der besser ist.
Aus diesem Gedanken des neuen Bundes heraus wird die gesamte Diskussion in Kapiteln 9 bis 10 fließen.
Der neue Bund ist besser, und seine Opfer und Riten sind besser.
Dieser neue Bund wird direkt mit der prophetischen Hoffnung in \bibverse{Jeremiah}{31}{} verbunden, das der Autor wortreich zitiert.
Dieser neue Bund wird `das Gesetz in ihren Sinn geben'\bibfoot{Hebrews}{8}{10}, und verbunden mit dem Schluss des neuen Bundes wird Gott `ihrer Sünden nicht mehr gedenken.'\bibfoot{Hebrews}{8}{12}.

Kapitel 9 bis 10 füllen jetzt die Details dieses Vergleiches aus. Wenn Jesus Hoher Priester eines neuen Bundes ist, was sind die Opfer und Riten dieses Bundes? Wie funktioniert er, und was bewirkt er?
Diese Kapitel werden wir jetzt Punkt für Punkt durchgehen und insbesondere sehen, wie Jesus als Opfer beschrieben wird.
Verse 11 bis 12 beginnen wie folgt:
\begin{quote}
Christus aber ist gekommen als Hoher Priester der zukünftigen Güter und ist durch das größere und vollkommenere Zelt –
das nicht mit Händen gemacht, das heißt, nicht von dieser Schöpfung ist – und nicht mit Blut von Böcken und Kälbern,
sondern mit seinem eigenen Blut ein für alle Mal in das Heiligtum hineingegangen und hat \textlangle uns\textrangle\ eine ewige Erlösung erworben
\attrib{\bibverse{Hebrews}{9}{11-12}, \Cite{elb}}
\end{quote}
Jesus ist besser als der alte Bund, und dieses besser beginnt damit, dass sein Blut besser ist als das von Tieren.
Wir erinnern uns an \bibverse{Leviticus}{17}{11} und die Diskussionen darüber, dass das Blut der Opfertiere oft der primäre Grund ist, sie zu Opfern.
Jesus Blut ist sein Leben, mehr als sein Tod.
Und mit seinem Leben (rituell angewandt, ausgeschüttet am Kreuz) tritt Jesus in das himmlische Heiligtum hinein.
Jesus \textit{ist} einfach von Natur aus der einzige Mensch, der aus sich selbst hinaus in das Heiligtum herauf steigen kann.\footnote{
Athanasius in \Cite{de-incarnatione-verbi} ist besonders fokussiert darauf, dass Jesus durch seine Inkarnation die göttliche Natur mit der menschlichen teilt,
was sehr gut mit der Feststellung harmonisiert, dass Jesus durch seine Natur das vollständige Leben ist, dass den Tod überwindet.
Wir werden in unserer Diskussion über die Kirchenväter dazu zurückkehren: Soteriologie ist Christologie.}

\begin{quote}
Denn wenn das Blut von Böcken und Stieren und die Asche einer jungen Kuh, auf die Unreinen gesprengt, zur Reinheit des Fleisches heiligt,
wie viel mehr wird das Blut des Christus, der sich selbst durch den ewigen Geist \textlangle als Opfer\textrangle\ ohne Fehler Gott dargebracht hat,
euer Gewissen zu reinigen von toten Werken, damit ihr dem lebendigen Gott dient!
\attrib{\bibverse{Hebrews}{9}{13-14}, \Cite{elb}}
\end{quote}
Hier ist das erste Ritual, das Jesus durch sein besseres Blut möglich macht: Jesus wäscht uns rein und heiligt uns.
Wir haben hier zwei Funktionen mit zwei verschiedenen Typen im Alten Bund:
(a) die körperliche Reinigung durch die Asche einer jugen Kuh bezieht sich auf \bibverse{Numbers}{19}{} und das Reinigungsritual für Leichen-Unreinheit.
So wie dieses physische Reinigungsmittel die rituelle Unreinheit von einem Menschen entfernt, der eine Leiche berührt hat, wäscht Jesu Blut unsere ``toten Werke''
ab.
(b) Das Blut von Böcken und Stieren auf die Unreinen gesprengt bezieht sich wohl auf die Anwendung von Blut beim Bundesschluss oder bei der Priesterweihe.\footnote{
	Die Blutanwendungen bei der Reinigung von Aussatz ist unwahrscheinlich, da das Blut dort auf den ehemals Aussätzigen `gegeben' wird, nicht `gesprengt'.
	Vergleiche \bibverse{Leviticus}{13}{}.
}
In \ref{metaphysische-transformation} haben wir bereits untersucht, was diese Besprengung bewirkt: eine Transformation hin zu einer größeren Heiligkeit.
Genau so ist es Jesu Blut, dass die Gläubigen reinigt und transformiert: es wäscht unsere toten Werke ab (wie die Asche einer jungen Kuh die Leichen-Unreinheit abwäscht),
und es macht uns heilig, wie das Blut gesprengt auf die werdenden Priester sie für ihre Aufgabe in Gottes direkter Gegenwart heiligt.
Das dekontaminierende Opfer hier nicht gemeint sein können sehen wir sofort daran, dass ihr Blut niemals auf Menschen angewandt wurde.
\footnote{Man könnte einwenden, dass προσήνεγκεν hier Jesus als Opfer darstellen muss und die Bedeutung als Reinigungsmittel ausschließt.
	Doch obwohl `opfern' eine natürlich Bedeutung ist, taucht προσφέρω in LXX genau so für nicht-sakrifiziale (in der Tat: nicht-rituelle) Handlungen auf.
	Die Einfügung \textlangle als Opfer\textrangle\ der \Cite{elb} ist hier nicht zwingend nötig, und dieses Verb kann beide Handlungen einschließen.
}
Jesu Blut ist das perfekte unausschöpfliche Leben par excellence - das beste Reinigungsmittel der Welt.
Und so reinigt es unser Gewissen von toten Werken, wo die Reinigung im levitischen System nur den Körper von der Unreinheit des Todes reinigen konnte.
Wir dienen einem lebendigen Gott, reingewaschen von unseren toten Werken und wir tun es durch die Waschung in Jesu göttlichem Leben -
hier sehen wir sofort die Antitypen der beiden Zentralen Punkte der prophetischen Hoffnung:
Waschung durch Gott, und Befähigung ihm zu dienen\footnote{
Denke an \bibverse{Jeremiah}{31}{34}, das der Autor gerade erst zitiert hat.
}.
Mit diesem neuen Bund schließt auch Verse 15 direkt an:
\begin{quote}
	Und darum ist er Mittler eines neuen Bundes, damit, da der Tod geschehen ist zur Erlösung von den Übertretungen unter dem ersten Bund, die Berufenen die Verheißung des ewigen Erbes empfangen.
	\attrib{\bibverse{Hebrews}{9}{15}, \Cite{elb}}
\end{quote}
Jesu Tod bringt `Erlösung der Übertretungen unter dem ersten Bund', und nicht nur das - er ermöglicht `den Berufenen die Verheißung des ewigen Erbes' zu empfangen.
Beide diese Konzepte sind direkt mit dem neuen Bundesschluss verbunden:
Der neue Bund ist es, der Vergebung der Sünden mit sich bringt nachdem unsere Natur und Sünde gewaschen wurde, genau wie Jeremiah vorher gesagt hatte.
Vergebung und Erlösung ist aus Gottes Gnade allein, durch den neuen Bund versiegelt im Opfer des Messias, für alle die an diesem Bund teilhaben.
Aber warum war es nötig, dass Jesus gestorben ist?
Hätte der neue Bund nicht ohne ein Opfer besiegelt werden können?
Zu dieser Frage wendet sich der Autor in den nächsten Versen:
\begin{quote}
	Denn wo ein Testament ist, da muss notwendig der Tod dessen eintreten, der das Testament gemacht hat. Denn ein Testament ist gültig, wenn der Tod eingetreten ist, weil es niemals Kraft hat, solange der lebt, der das Testament gemacht hat.
	\attrib{\bibverse{Hebrews}{9}{16-17}, \Cite{elb}}
\end{quote}
Interessant ist hier vor allem, dass der Autor mit der konsistenten Metapher- und Typ-Antityp-Beziehung zum levitischen System bricht und plötzlich eine legale Metapher verwendet, um zu zeigen, warum Jesus Tod notwendig war.
Dass er das tut ist aber sehr verständlich wenn wir beachten, dass Opfer im levitischen System so explizit als Nicht-Tode formuliert wurden (wie in Kapitel \ref{opfer-als-nicht-tod} besprochen).
Jesu Opfer war nötig, damit sein Blut zu unserer Waschung zur Verfügung steht.
Sein tatsächlicher Tod war auch nötig - aber um das zu begründen bedient sich der Autor bewusst einer Metapher, die nicht das Opfer direkt mit dem Tod in Verbindung bringt.
Zu der Notwendigkeit des Blutes kehren wir prompt zurück:
\begin{quote}
Daher ist auch der erste ⟨Bund⟩ nicht ohne Blut eingeweiht worden.
Denn als jedes Gebot nach dem Gesetz von Mose dem ganzen Volk mitgeteilt war, nahm er das Blut der Kälber und Böcke mit Wasser und Purpurwolle und Ysop und besprengte sowohl das Buch selbst als auch das ganze Volk
und sprach: »Dies ist das Blut des Bundes, den Gott für euch geboten hat.«
Aber auch das Zelt und alle Gefäße des Dienstes besprengte er ebenso mit dem Blut;
	\attrib{\bibverse{Hebrews}{9}{18-21}, \Cite{elb}}
\end{quote}
Wieder und wieder dreht sich hier alles um Jesu Blut als Blut des Bundes, das selbst im Heiligtum angebracht wurde zur Ratifizierung des Bundes.
Es wäscht, aber vor allem heiligt es, wie wir in Kapitel \ref{metaphysische-transformation} gesehen haben.
Im folgenden möchte der Autor Jesu Werk für die Dekontamination des himmlischen Heiligtums erläutern, und um den Bogen vom Blut des Bundes zum Blut der Dekontamination zu schlagen steht Vers 22:
\begin{quote}
	und fast alle Dinge werden mit Blut gereinigt nach dem Gesetz, und ohne Blutvergießen gibt es keine Vergebung.
	\attrib{\bibverse{Hebrews}{9}{22}, \Cite{elb}}
\end{quote}
Der erste Satz ist leicht verständlich: so wie alle Dinge mit Blut geheiligt werden, werden sie auch dekontaminiert.
Der zweite Teilsatz ist der komplizierteste Vers in diesen zwei Kapiteln des Hebräerbriefes, einfach weil er - naiv gelesen - falsch ist.
Wir haben bereits im Detail beobachtet, dass Vergebung unabhängig von Blut und Opfern geschehen kann (in Kapitel \ref{chap-fiat}).
Außerdem wissen wir, dass das Resultat der Säuberungsopfer die in v22a offensichtlich im Fokus stehen nicht Vergebung, sondern Dekontamination ist (Kapitel \ref{kipper}).
Nicht nur das, sondern selbst wenn Säuberungsopfer Sünden vergeben würden, wird nicht bei allen Säuberungsopfern Blut vergossen.
Hier sind diverse Möglichkeiten, was `ohne Blutvergießen gibt es keine Vergebung' bedeuten kann:
\begin{description}
	\item[Gott kann keine Sünden vergeben, ohne dass Blut fließt] Diese Aussage ist wie schon diskutiert einfach falsch.
	\item[Freilassung durch den neuen Bund] Diese Idee passt zwar zur bisherigen Diskussion in Kapitel 9, aber sie unterbricht den Umschwung der Argumentation von Bundesschluss zu Jesus als Säuberungsopfer und ist daher unwahrscheinlich.
	\item[Freilassung von der Notwendigkeit, Säuberungsopfer zu bringen] Diese Aussage ist mit Sicherheit richtig, erscheint aber recht kompliziert.
	\item[Freilassung von der Dekontamination an sich] Diese Option muss abgelehnt werden weil sie verlangt, dass ἄφεσις mit \kipper\ gleichzusetzen ist - diese Gleichsetzung ist dadurch unwahrscheinlich, dass die LXX niemals ἄφεσις (oder seine verbale Form) für kipper verwendet, sondern nur für נִסְלַח (nislaḥ, vergeben) in diesen Kontexten.
		Wenn der Kontext für die Verwendung von ἄφεσις im levitischen System zu finden ist, dann kann dieses Wort nur schwer Dekontamination bedeuten.
	\item[Vergebung von Sünden für die es Säuberungsopfer gab] Auch diese Aussage erscheint kompliziert, aber zumindest fließt sie gut mit dem Argument.
		(a) Es geht tatsächlich weiter um Säuberungsopfer
		(b) das Blutvergießen wird konsistent mit v22a und v23 als Blut eines dekontaminierenden Opfers gelesen
		(c) diese Aussage ist explizit eine Vorschrift aus dem Gesetz, das in v22 angebracht wird und
		(d) bis auf das Problem der Säuberungsopfer für Arme, die ohne Blut ablaufen, ist die Aussage tatsächlich richtig und wird nicht hier das erste mal so formuliert.\footnote{
		Gott wird Sünden für die es Säuberungsopfer gibt nur vergeben, wenn diese Opfer auch gebracht werden.
	Diese Aussage ist so unproblematisch, dass sie in einem Halbsatz ausgesprochen werden kann ohne vom eigentlichen Argument abzulenken.}
\end{description}
Auch die letzte Option erscheint etwas kompliziert, aber es ist meiner Ansicht nach die beste Interpretation dieses sehr schwieriegen Verses.

Die Argumentation des Briefes schreitet jetzt fort:
\begin{quote}
\textlangle Es ist\textrangle\ nun nötig, dass die Abbilder der in den Himmeln \textlangle befindliche Dinge\textrangle\ hierdurch gereinigt werden, die himmlischen Dinge selbst aber durch bessere Schlachtopfer als diese.
Denn Christus ist nicht hineingegangen in ein mit Händen gemachtes Heiligtum, ein Abbild des wahren \textlangle Heiligtums\textrangle\, sondern in den Himmel selbst,
um jetzt vor dem Angesicht Gottes für uns zu erscheinen, auch nicht, um sich selbst oftmals zu opfern, wie der Hohe Priester alljährlich mit fremdem Blut in das Heiligtum hineingeht
– sonst hätte er oftmals leiden müssen von Grundlegung der Welt an –; jetzt aber ist er einmal in der Vollendung der Zeitalter offenbar geworden, um durch sein Opfer die Sünde aufzuheben.
	\attrib{\bibverse{Hebrews}{9}{23-26}, \Cite{elb}}
\end{quote}
So wie ein irdisches Säuberungsopfer das irdische Heiligtum reinigt, so reinigt Jesu perfektes Blut das himmlische Heiligtum.
Das ist eine Facette der `Versöhnung', die in Christus passiert: alle Unreinheit, die durch Sünde\footnote{und ggf. rituelle Unreinheit, aber ob rituelle Reinheit für den Autor noch eine Rolle spielt ist sehr fragwürdig} auf das himmlische Heiligtum gekommen ist wird gereinigt.
Nur so ist es möglich, dass das himmlische Heiligtum seine Funktion weiterhin erfüllen kann.
Das ist `Versöhnung' durch sein Blut - das Heiligtum ist gereinigt und wir haben Zugang zum Vater, durch Jesu Opfertod am Kreuz.
Das Problem, dass der Hebräerbrief hier löst ist `das himmlische Heiligtum ist verunreinigt' - und die Lösung ist Jesu Opfertod für unsere Sünden.
Substitionelle Bestrafung des perfekten Hohe Priesters des neuen Bundes, der sich selbst als Opfer darbringt ist in dieser Passage nicht im Ansatz indikiert.
Wenn Jesu Tod als stellvertretende Strafe zur Vergebung von Sünden das Ziel dieser Passage wäre, müsste der Autor gleich mit mehreren Tabus des levitischen Systems brechen (Opfer sind keine Strafe; Opfer sind kein stellvertretender Tod; Opfer sind nicht einmal Tod; Opfer haben eine sehr klar definierte andere Funktion), aber das ist nicht was wir sehen.
Der Autor kennt sein Levitikus, und er verwendet das vielschichtige rituelle System der Torah mit großer Vorsicht und Klarheit.
Der Messias und sein Tod als Opfer für unsere Sünden erfüllt alle rituelle Funktionen des levitischen Systems.
Wir sind reingewaschen, wir sind geheiligt als Teilhaber an einen neuen Bund, das Heiligtum ist gereinigt.
Diese bereits etablierten Funktionen werden jetzt durch den Rest der Passage noch weiter ausformuliert, um weiter klar zu machen dass Jesu Opfer größer ist als es die levitischen jemals sein konnten.

Verse 27 und 28 bilden eine Vorausschau auf Jesu zweites Kommen und das endgültige Gericht, und zum Beginn von Kapitel 10 fahren wir mit unserem Thema - Jesus und das levitische System - fort.
\begin{quote}
Denn da das Gesetz einen Schatten der zukünftigen Güter, nicht der Dinge Ebenbild selbst hat, so kann es niemals mit denselben Schlachtopfern,
die sie alljährlich darbringen, die Hinzunahenden für immer vollkommen machen.
Denn würde sonst nicht ihre Darbringung aufgehört haben, weil die den Gottesdienst Übenden, einmal gereinigt, kein Sündenbewusstsein mehr gehabt hätten?
Doch in jenen \textlangle Opfern\textrangle\ ist alljährlich ein Erinnern an die Sünden;
denn unmöglich kann Blut von Stieren und Böcken Sünden wegnehmen.
	\attrib{\bibverse{Hebrews}{10}{1-4}, \Cite{elb}}
\end{quote}
Diese Aussagen sollten uns zum jetzigen Zeitpunkt sehr bekannt vorkommen: Natürlich können die levitischen Opfer keine Sünden wegnehmen - nie war das ihre Funktion.
Dass das eigentliche Problem, nämlich die Sünde an sich, von ihnen nie gelöst wurde sehen wir daran, dass diese Opfer immer wieder dargebracht werden müssen.
Jetzt aber hat Christus ein Opfer gebracht, nicht bloß um das Heiligtum zu reinigen und nicht bloß um Vergebung auszusprechen - sondern um mit diesem Opfer eine wirkliche Transformation seines Bundesvolkes zu bewirken:
\begin{quote}
Darum spricht er, als er in die Welt kommt: »Schlachtopfer und Opfergabe hast du nicht gewollt, einen Leib aber hast du mir bereitet;
an Brandopfern und Sündopfern hast du kein Wohlgefallen gefunden.
Da sprach ich: Siehe, ich komme – in der Buchrolle steht von mir geschrieben –, um deinen Willen, Gott, zu tun.«
Vorher sagt er: »Schlachtopfer und Opfergaben und Brandopfer und Sündopfer hast du nicht gewollt, auch kein Wohlgefallen daran gefunden« –
die doch nach dem Gesetz dargebracht werden –;
dann sprach er: »Siehe, ich komme, um deinen Willen zu tun« – er nimmt das Erste weg, um das Zweite aufzurichten.
In diesem Willen sind wir geheiligt durch das ein für alle Mal geschehene Opfer des Leibes Jesu Christi.
	\attrib{\bibverse{Hebrews}{10}{5-10}, \Cite{elb}}
\end{quote}
Denn das Opfer Jesu bewirkt etwas größeres, als all die Opfer des alten Bundes:
Jesu Opfer `heiligt' uns in Analogie dazu, wie die Opfer des Bundesschlusses das Volk geheiligt haben, aber auf eine viel bessere und größere Art und Weise.\footnote{
\bibverse{Psalms}{40}{}, aus dem hier zitiert wird, enthält eine der beinahe üblichen Aussprüche darüber, dass Gott das Herz wichtiger ist als die Opfer an sich.
Außerdem versteht der Autor des Hebräerbriefes diesen Psalm scheinbar messianisch, und stellt in seinem Midrasch Jesu Wirken und speziell sein freiwilliges Opfer als das da, was Gott wirklich wollte:
Den perfekten Gott-Mensch, dessen Leben und Tod Gott ein wohlgefälliges Opfer ist.}
An den nächsten Versen können wir sehen, dass der Fokus des Autors hier tatsächlich konsistent auf Jesu Opfer als Antitypos des Bundesschlusses -
und insbesondere der metaphysischen Transformation des Bundesvolkes - liegt.
Insbesondere ist also der kurze Ausflug in eine Diskussion über Jesus als Säuberungsopfer bereits beendet:
\begin{quote}
Und jeder Priester steht täglich da, verrichtet den Dienst und bringt oft dieselben Schlachtopfer dar, die niemals Sünden hinwegnehmen können.
Dieser aber hat ein Schlachtopfer für Sünden dargebracht und sich für immer gesetzt zur Rechten Gottes.
Fortan wartet er, bis seine Feinde hingelegt sind als Schemel seiner Füße.
Denn mit einem Opfer hat er die, die geheiligt werden, für immer vollkommen gemacht.
⟨Das⟩ bezeugt uns aber auch der Heilige Geist; denn nachdem er gesagt hat:
»Dies ist der Bund, den ich für sie errichten werde nach jenen Tagen, spricht der Herr, ich werde meine Gesetze in ihre Herzen geben und sie auch in ihren Sinn schreiben«;
und: »Ihrer Sünden und ihrer Gesetzlosigkeiten werde ich nicht mehr gedenken.«
	\attrib{\bibverse{Hebrews}{10}{11-17}, \Cite{elb}}
\end{quote}
Wir sind vielleicht geneigt zu lesen ``Dieser aber hat ein Schlachtopfer für Sünden dargebracht'' und daraus zu lesen, dass `Sein stellvertretendes Sühneopfer unsere Schuld auf sich nimmt', aber eine solche Interpretation ist mit einem klaren Verständnis der Argumentation und der Hintergründe im levitischen System einfach unmöglich.
Um es nocheinmal zu sagen:
Kein Opfer hat jemals `stellvertretend Sünden vergeben/bestraft',
der Autor hat nie argumentiert, dass sich dieser Zustand jetzt ändert,
und wir dürfen den Text nicht in unser falsches Verständnis des levitischen Systems zwingen.
Tatsächlich ist der Text sehr eindeutig darin, welche Opfer hier angesprochen werden und in welcher Art und Weise genau diese Opfer `für Sünden' sind.
Bereits in Vers 10 war das Ziel des Opfers `Heiligung'; in Vers 12 ist es `für Sünden' und in Vers 14 zur `Vervollkommnung'.
Direkt danach, und direkt als Begründung verbunden wird wiederum die Hoffnung des neuen Bundes aus \bibverse{Jeremiah}{31}{} zitiert.
Jesu Opfer in Versen 10 bis 14 ist ganz eindeutig ein Opfer des Bundesschlusses, kein Säuberungsopfer.
Jesu Blut, das am Kreuz für unsere Sünden vergossen wurde läutet den neuen Bund ein, auf den die Propheten gehofft haben.
Es heiligt uns und bringt uns zur Vollkommenheit, indem es uns transformiert.
Ja, unsere Sünden sind uns vergeben - denn so einen Gott haben wir, der Sünder zu seinem Bundesvolk macht und ihnen vergibt.
Aber das Opfer des fleischgewordenen Gottes bewirkt viel mehr als das.
Das perfekte Blut des inkarnierten Wortes, der nicht anders kann als zu leben ist es, durch das wir \textit{leben}, durch den wir \textit{heilig sind} und durch das \textit{die Ursache für die Sünde} zusammen mit der Strafe für diese Sünde ausgelöscht wird: unsere gefallene, korrupte Natur.
Dazu war Christi Tod nötig: um den neuen Bund einzuleiten - und genau das hat der Autor explizit in \bibverse{Hebrews}{9}{16-17} argumentiert.

\begin{quote}
Wo aber Vergebung dieser \textlangle Sünden\textrangle\ ist, gibt es kein Opfer für Sünde mehr.
	\attrib{\bibverse{Hebrews}{10}{18}, \Cite{elb}}
\end{quote}
Hier begegnet uns eine besondere Feinheit, die in der deutschen Übersetzung leider schwer zu transportieren ist:
Der technische Term für die Säuberungsopfer in der LXX ist περὶ ἁμαρτίας (peri hamartias, für Sünden), und diese Konstruktion ist uns in den Zitaten aus Psalm 40 begegnet, wo die Elberfelder Bibel `Sündopfer' übersetzt hat.
Interessanterweise hat der Autor in Vers 12 `ein Schlachtopfer für Sünden' hier mit ὑπὲρ ἁμαρτιῶν προσενέγκας θυσίαν (hyper hamartion prosenegkas thysian) wiedergegeben - einem anderen grichischen Wort das hier im Deutschen legitim `für' übersetzt werden kann.
In Vers 18 kehrt der Autor jetzt zu περὶ ἁμαρτίας zurück - dem griechischen Begriff, der explizit mit Säuberungsopfern verbunden ist.
Wo ein neuer Bund geschaffen ist, um die Sünden zu vergeben und abzuwaschen gibt jetzt keine anderen Säuberungsopfer mehr, um sich um die Sünden zu kümmern.\footnote{
Selbst die Sünden, für die vorher Säuberungsopfer vorgeschrieben waren - alles ist mit einem finalen, perfekten Opfer beseitigt.}
Der neue Bund und das neue Opfer ist so viel besser, dass ``er den ersten für veraltet erklärt''\bibfoot{Hebrews}{8}{13}.

In den folgenden Versen wird das theoretisch Erklärte für die Zuhörer praktisch:
\begin{quote}
Da wir nun, Brüder, durch das Blut Jesu Freimütigkeit haben zum Eintritt in das Heiligtum,
den er uns eröffnet hat als einen neuen und lebendigen Weg durch den Vorhang – das ist durch sein Fleisch –,
und einen großen Priester über das Haus Gottes,
so lasst uns hinzutreten mit wahrhaftigem Herzen in voller Gewissheit des Glaubens,
die Herzen besprengt ⟨und damit gereinigt⟩ vom bösen Gewissen und den Leib gewaschen mit reinem Wasser.
	\attrib{\bibverse{Hebrews}{10}{19-22}, \Cite{elb}}
\end{quote}
Jesu Opfer hat nicht nur den neuen Bund eingeleitet, Sünden vergeben und uns rein gemacht - er eröffnet uns auch den Zugang durch den Vorhang, in das Allerheiligste.
Mit unserem Verständnis des levitischen Systems ist auch diese Aussage direkt aus dem vorher gesagten logisch abzuleiten:
Wir sind rein gewaschen durch Jesu perfektes reinigendes Blut,
wir sind besprengt mit dem Blut des Bundes - und damit metaphysisch transformiert hin zu einem heiligen Volk Gottes.
So sind wir als Menschen also vorbereitet, das Heiligtum zu betreten.
Nicht nur wir, sondern auch das himmlische Heiligtum wurde vorbereitet durch Jesu Blut als Säuberungsopfer, mit dem das Heiligtum gereinigt wurde.
Alles ist bereit, alles ist in kultischer Ordnung und jetzt ist der Weg frei;
durch Jesu Leben, durch sein Blut, mit ihm als unserem Fürsprecher vollständig mit Gott versöhnt und ohne Angst in seiner Gegenwart.

Es folgt eine Ermahnung, an dem Bekentniss dieses Christus festzuhalten, gefolgt von dieser weiteren Warnung, dass Jesu endgültiges Opfer nicht verachtet werden darf:
\begin{quote}
Denn wenn wir mutwillig sündigen, nachdem wir die Erkenntnis der Wahrheit empfangen haben, bleibt kein Schlachtopfer für Sünden mehr übrig,
sondern ein furchtbares Erwarten des Gerichts und der Eifer eines Feuers, das die Widersacher verzehren wird.
Hat jemand das Gesetz Moses verworfen, stirbt er ohne Barmherzigkeit auf zwei oder drei Zeugen hin.
Wie viel schlimmere Strafe, meint ihr, wird der verdienen, der den Sohn Gottes mit Füßen getreten und das Blut des Bundes,
durch das er geheiligt wurde, für gemein erachtet und den Geist der Gnade geschmäht hat?
	\attrib{\bibverse{Hebrews}{10}{26-29}, \Cite{elb}}
\end{quote}
Hier zeigt sich wieder, was das Zentrum Jesu Wirkens ist: der neue Bund.
Denn seinem Bundesvolk in Christus hat der Vater völlig vergeben, sie angenommen und geheiligt.
Es gibt keine andere Möglichkeit, dem kommenden Gericht zu entkommen, als Teil dieses Bundes zu sein -
kein anderes Opfer ist in der Lage das zu ersetzen, was Christus auf viel bessere Art und Weise bewirkt hat.


Der Brief an die Hebräer ist für unser Verständnis des Opfers am Kreuz eine wahre Fundgrube an komplizierten Bezügen auf das levitische System und die prophetische Hoffnung
des neuen Bundes -
Bezüge, die plötzlich ein wunderschönes und klares Bild ergeben,
nachdem wir in den ersten Kapiteln die Farben und die Pinsel einzeln under die Lupe genommen haben.
Das Kreuz - das freiwillige Opfer des inkarnierten Wortes zur Vergebung unserer Sünden - ist ein Höhepunkt der Geschichte eines liebenden Gottes,
der seinem Bundesvolk großzügig und unverdient vergibt; der sein göttliches Blut des perfekten Lebens vergießt um uns rein zu waschen und zu heilen aus unserer korrupten Natur;
der seinen Sohn in den Tod gibt, damit sein Blut das himmlische Heiligtum reinigen kann, sodass wir unter der unaufhörlichen Fürsprache des immer lebenden Sohnes,
gewaschen im Heiligen Geist, und ohne Furch vor den Vater treten können.

\chapter{Gerechtigkeit und das letzte Gericht}
Ein wichtiger Bestandteil der klassischen protestantischen (und insbesondere reformierten) Soteriologie ist eine bestimmte
Ansicht darüber, auf welcher Basis die Gläubigen durch das letzte Gericht hindurch bestehen werden.
Die Kurzfassung ist diese:
Am Ende der Zeit wird jede Person vor Gottes Gericht erscheinen. Die nicht durch Jesus erlösten werden wegen ihrer Sünden zum Tode verurteilt.
Für die Erlösten sieht Gott statt ihrer eigenen Gerechtigkeit auf die Perfektion Jesu, die uns angezogen wurde und verdammt uns auf dieser Grundlage nicht.

Zu dieser Ansicht kommt das Verständnis hinzu, dass Gott jede Sünde bestrafen muss, und um uns nicht zu vernichtet stattdessen Jesus die Strafe auf sich nimmt
- diese zweite Seite der Medaillie ist dann PSA.
In diesem Artikel ist unsere eigentliche Frage die zweite: hat der Vater den Sohn für unsere Schuld bestraft?
Wenn wir diese Frage mit Nein beantworten, wie ich es tue, ist es aber nötig ein weiteres loses Ende zu erklären:
Was bedeutet es dann, dass wir durch Jesu Tod gerechtfertigt werden?
Dieser Frage wenden wir uns in diesem Kapitel vor. Gleichzeitig können wir beim beantworten dieser Frage
natürlich über einige Bibelstellen sprechen, die gerne in Verteidigungen von PSA verwendet werden - aber eine bessere Interpretation haben.

Hier ist meine Sicht vorweggenommen:
(a) Jeder Mensch wird im letzten Gericht für seine tatsächliche Gerechtigkeit nach seinen Werken gerichtet.
Die Gerechten werden leben, die Ungerechten werden vernichtet.
(b) Die griechischen Wörter hinter `Gerechtfertigung' und seiner verbalen und adjektivischen Form sind nicht ausschließlich forensisch, sondern
können ontologisch fungieren\footnote{
	Das heißt, `Gerechtigkeit', `gerecht machen', ... sind in manchen Kontexten valide Übersetzungen.
}
(c) Wenn Paulus diese Wörter im Römerbrief verwendet, sind sie besser in diesem ontologischen Charakter zu verstehen:
durch Glauben an Jesus werden wir gerecht gemacht, nicht bloß unabhängig von unseren Werken gerecht gesprochen.
Diese Konzepte so zu verstehen, und insbesondere den Römerbrief so zu lesen, ist die beste Möglichkeit, zwei Wahrheiten zu verstehen, die
häufig gegensätzlich erscheinen: Wir sind gerichtet nach Werken, aber gerechtfertigt durch Glauben.\footnote{
	Ich habe nicht genug Platz, um alle Verse zu besprechen, die in dieser Diskussion interessant wären.
	Ein gutes Beispiel ist \bibverse{Ephesians}{2}{8-9} (beachte zu diesem speziellen Vers: was meint Paulus hier mit `gerettet'?).
	Da mein primäres Ziel in diesem Kapitel ist, Paulus' Verständnis von Gerechtfertigung in Relation zu Jesu Tod und Leben zu verstehen
	bin ich hier nicht erschöpfend.
}

\section{\label{sec:gericht-nach-werken}Gericht nach Werken}
Die jüdische Literatur in der Zeit des zweiten Tempels ist in vielen speziellen Fragen recht divers.
Doch zumindest in dieser Frage gibt es keine Ausreißer: jeder Mensch wird am Ende nach seinen Werken gerichtet.
Wer gerecht gelebt hat wird leben, und wer ungerecht war wird vernichtet.
Wenn eine Person vor diesem Gericht besteht, dann weil sie tatsächlich gut gelebt hat.
In diesem Gericht sieht Gott die Person nicht an, und wird niemanden belohnen, der diese Belohnung nicht verdient hat.
Wie absolut und ausnahmslos dieses Verständnis ist hat VanLandingham in seinem Buch \Cite{judgement-and-justification-in-early-judaism-and-the-apostle-paul} gezeigt.
Hier gebe ich nur eine kurze Auswahl dieser Stellen wieder:

\begin{quote}
	Zerrissen wird die Erde,
und alles auf der Erde wird vergehen,
und alles wird gerichtet werden.
Doch mit den Frommen schließt Er Frieden
und schützt die Auserwählten
und Gnade waltet über ihnen.
Sie werden alle Gottes Eigentum
und sind im Glücke und gesegnet.
Er selber unterstützt sie alle
und Gottes Licht wird ihnen scheinen.
Er selber schließt mit ihnen Frieden.
Fürwahr! Er kommt mit Tausenden von Heiligen,
um über alle das Gericht zu halten
und alle Übeltäter zu vernichten
und alles Fleisch zurechtzuweisen
der schlimmen Taten wegen,
die sie so frevlerisch begingen,
sowie der kühnen Worte halber,
die gegen Ihn die Sünder frevelhaft gesprochen.\footnote{
	In \Cite[5]{aeth-hen} lesen wir weiterhin, dass die Gottes Gesetz nicht beachten `keinen Frieden finden' (4) werden, während für die Gerechten
	während oder nach dem Gericht gilt: `Und Weisheit wird verliehen den Auserwählten; sie leben all und sündigen nicht mehr'.(8)
	Wir werden später sehen, wie ähnlich diese Idee Paulus Argument im Römerbrief kommt, bloß dass Paulus das Ausgießen von Gottes Gerechtigkeit auf den Christus
	bezieht, nicht auf die Zeit nach dem letzten Gericht.
}
\attrib{\Cite[1:7-9]{aeth-hen}}
\end{quote}

\begin{quote}
Reflect on all these matters and seek from him that he may support your
counsel and keep far from you the evil schemings and the counsel of Belial,
so that at the end of time, you may rejoice in finding that some of our words
are true. And it shall be reckoned to you as justice when you do what is
upright and good before him, for your good and that of Israel.\footnote{
	Für weiteres Material aus Qumran, beachte z.B. die Zwei-Wege Kapitel in 1QS 3-4.
}
\attrib{\Cite[4QMMTe Frag. 14-17 Col II]{tdss-se}}
\end{quote}

\begin{quote}
	Wer rechtschaffen handelt, erwirbt sich Leben beim Herrn,
  und wer Unrecht thut, verwirkt selbst sein Leben in Verderben
	denn des Herrn Gerichte sind gerecht gegen Person und Haus.\footnote{
		Vergleiche auch \Cite[9:11]{pss-sol} (Des Herrn ist das Erbarmen über das Haus Israel immer und ewig.) mit \Cite[10:8]{pss-sol} (Des Herrn [Werk] ist die Erlösung über das Haus Israel, zur ewigen Freude.).
		Wie VanLandingham beschreibt ist hier die Gottes Erbarmen (seine Gnade) gleichbedeutend mit der Errettung als solche, nicht mit einer bestimmten Art und Weise auf die 
		die Erlösten anders als die Ungerechten im letzten Gericht behandelt werden. Vergleiche \bibverse{Isaiah}{56}{1} LXX für einen ähnlichen Parallelismus.
	}
	\attrib{\Cite[9:5]{pss-sol}}
\end{quote}

\begin{quote}
	Ich antwortete und sprach: Herr Gott, du hast ja in deinem Gesetze bestimmt,
	nur die Gerechten würden dies Erbteil bekommen, aber die Gottlosen sollten ins Verderben gehen.
	\attrib{\Cite[7:17]{fourth-ezra}}
\end{quote}

\begin{quote}
	Denn sieh, es kommen Tage;
da öffnet man die Bücher,
worin die Sünden aller Missetäter aufgeschrieben,
sowie die Kammern,
wo die Gerechtigkeit all derer aufgespeichert ist,
die in der Schöpfung recht gehandelt.
\attrib{\Cite[24:1]{second-baruch}}
\end{quote}

Viele dieser Texte basieren auf der Deuteronomischen Formel\bibfoot{Deuteronomy}{28}{}:
Wenn Israel Gott gehorcht werden sie gesegnet sein und erhaben sein - hier sehr praktisch als Nation auf dieser Welt.
Wenn sie aber nicht gehorchen, wird die Nation verflucht sein und ins Exil gehen.
Genau das geschieht, und Israel geht ins Exil - nicht weil Gott zu schwach ist, sondern weil das Volk gesündigt hat und Gott Gericht hält.
Selbst als Juda in sein Land zurückgekehrt ist sündigt das Volk weiter, und die Flüche des Bundes lasten weiter auf ihm.\footnote{
	\bibverse{Haggai}{1}{}, das Buch Maleachi, \bibverse{Zechariah}{11}{4-17}.
	In vielen anderen Stellen ringt Gott durch die Propheten um sein Volk, dass es nicht in die Wege ihrer Väter zurück fällt: \bibverse{Zechariah}{1}{1-6}.
	Und die Rückkehrer sahen ihre schwere Lage zumindest als Resultat der Sünde ihrer Väter: \bibverse{Nehemiah}{9}{32-37}, \bibverse{Ezra}{9}{5ff}.
	Wir haben in Kapitel \ref{chap:bundesfluch} bereits gesehen, dass die Juden zur Zeit des Neuen Testaments immer noch der Meinung waren, dass
	das Exil nicht völlig beendet ist.
}.

In dieser Situation sehen wir in Maleachi die Hoffnung, dass Gott an seinem Tag die gerechten und die gottlosen voneinander trennen wird\footnote{
	Vergleiche auch \bibverse{Isaiah}{65}{}, wo die Endzeitliche Hoffnung auf die Gerechten beschränkt ist und die Ungerechten nicht teilhaben.
	Ebenso \bibverse{Daniel}{12}{2-3}, in dem aber nur eine der Seiten (die Gerechten, die viele zur Gerechtigkeit führen) explizit beschrieben wird.
}:
\begin{quote}
	Denn siehe, der Tag kommt, der wie ein Ofen brennt. Da werden alle Frechen und alle, die gottlos handeln, Strohstoppeln sein. Und der kommende Tag wird sie verbrennen, spricht der HERR der Heerscharen, sodass er ihnen weder Wurzel noch Zweig übrig lässt.
	Aber euch, die ihr meinen Namen fürchtet, wird die Sonne der Gerechtigkeit aufgehen, und Heilung ist unter ihren Flügeln. Und ihr werdet hinausgehen und umherspringen wie Mastkälber.
	\attrib{\bibverse{Malachi}{3}{19-20}}
\end{quote}

Und auch im neuen Testament findet sich dieselbe Ansicht ohne große Diskussion\footnote{
	Und das sollten wir auch erwarten: Entweder der Jesus-Glauben reiht sich in das Verständnis aller anderen Juden der Zeit ein,
	oder wir erwarten eine klare Diskussion darüber, warum dieses Gericht doch anders aussehen wird.
}:
\begin{quote}
	Wenn aber der Sohn des Menschen kommen wird in seiner Herrlichkeit und alle Engel mit ihm, dann wird er auf seinem Thron der Herrlichkeit sitzen;
	und vor ihm werden versammelt werden alle Nationen, und er wird sie voneinander scheiden, wie der Hirte die Schafe von den Böcken scheidet.
	Und er wird die Schafe zu seiner Rechten stellen, die Böcke aber zur Linken.
	Dann wird der König zu denen zu seiner Rechten sagen: Kommt her, Gesegnete meines Vaters, erbt das Reich, das euch bereitet ist von Grundlegung der Welt an!
	Denn mich hungerte, und ihr gabt mir zu essen; mich dürstete, und ihr gabt mir zu trinken; ich war Fremdling, und ihr nahmt mich auf;
	nackt, und ihr bekleidetet mich; ich war krank, und ihr besuchtet mich; ich war im Gefängnis, und ihr kamt zu mir.
	Dann werden die Gerechten ihm antworten und sagen: Herr, wann sahen wir dich hungrig und speisten dich? Oder durstig und gaben dir zu trinken?
	Wann aber sahen wir dich als Fremdling und nahmen dich auf? Oder nackt und bekleideten dich?
	Wann aber sahen wir dich krank oder im Gefängnis und kamen zu dir?
	Und der König wird antworten und zu ihnen sagen: Wahrlich, ich sage euch, was ihr einem dieser meiner geringsten Brüder getan habt, habt ihr mir getan.
	Dann wird er auch zu denen zur Linken sagen: Geht von mir, Verfluchte, in das ewige Feuer, das bereitet ist dem Teufel und seinen Engeln!
	Denn mich hungerte, und ihr gabt mir nicht zu essen; mich dürstete, und ihr gabt mir nicht zu trinken;
	ich war Fremdling, und ihr nahmt mich nicht auf; nackt, und ihr bekleidetet mich nicht; krank und im Gefängnis, und ihr besuchtet mich nicht.
	Dann werden auch sie antworten und sagen: Herr, wann sahen wir dich hungrig oder durstig oder als Fremdling oder nackt oder krank oder im Gefängnis und haben dir nicht gedient?
	Dann wird er ihnen antworten und sagen: Wahrlich, ich sage euch, was ihr einem dieser Geringsten nicht getan habt, habt ihr auch mir nicht getan.
	Und diese werden hingehen zur ewigen Strafe, die Gerechten aber in das ewige Leben.
	\attrib{\bibverse{Matthew}{25}{31-46}, \Cite{elb}}
\end{quote}
Diese Gerichtszene auf Jesu Lippen trägt alle üblichen Charakterisika von Szenen über das letzte Gericht:
Gott - hier der Menschensohn - hält Gericht mit seinen Engeln. Die Menschen werden in zwei Gruppen eingeteilt: die Gerechten und die Ungerechten.
Diese Aufteilung der Menschen geschieht basierend auf Werken - hier speziell basierend auf der Nächstenliebe, die jeder gezeigt hat.
Die einzige Neuerung in dieser Beschreibung ist, dass der Menschensohn das Gericht hält, nicht der Vater selbst.

\begin{quote}
	Wundert euch darüber nicht, denn es kommt die Stunde, in der alle, die in den Gräbern sind, seine Stimme hören
	und hervorkommen werden; die das Gute getan haben zur Auferstehung des Lebens,
	die aber das Böse verübt haben zur Auferstehung des Gerichts.
	\attrib{\bibverse{John}{5}{28-29}, \Cite{elb}}
\end{quote}
Wieder gibt es ein Gericht über alle Menschen, die Entscheidung fällt basierend auf dem Guten oder Bösen, dass jeder getan hat, und das Ergebnis ist Leben oder Gericht.

\begin{quote}
	Denn wir müssen alle vor dem Richterstuhl Christi offenbar werden,
	damit jeder empfängt, was er durch den Leib \textlangle vollbracht\textrangle,
	dementsprechend, was er getan hat, es sei Gutes oder Böses.
	\attrib{\bibverse{IICorinthians}{5}{10}, \Cite{elb}}
\end{quote}
Hier sehen wir, dass Paulus sich wie alle anderen Juden seiner Zeit in diesen Konsens einreiht:
Es gibt ein endgültiges Gericht, dort wird jeder für seine Taten
gerichtet werden. Die einzige Inovation ist, dass Christus der Richter ist, nicht der Vater.
Am deutlichsten finden wir diese Wahrheit in Paulus Briefen im zweiten Kapitel des Römberbriefs.
Auch in dieser Stelle weicht Paulus nicht von der gewöhnlichen Ansicht der Juden zum ewigen Gericht ab: 
Gericht am Ende der Zeit, basierend auf Werken, mit dem Resultat von Leben oder Gottes endgültige Strafe.

\begin{quote}
	Nach deiner Störrigkeit und deinem unbußfertigen Herzen aber häufst du dir selbst Zorn auf für den Tag des Zorns und der Offenbarung des gerechten Gerichts Gottes,
	der einem jeden vergelten wird nach seinen Werken:
	denen, die mit Ausdauer in gutem Werk Herrlichkeit und Ehre und Unvergänglichkeit suchen, ewiges Leben;
	denen jedoch, die von Selbstsucht \textlangle bestimmt\textrangle und der Wahrheit ungehorsam sind, der Ungerechtigkeit aber gehorsam, Zorn und Grimm.
	Bedrängnis und Angst über die Seele jedes Menschen, der das Böse vollbringt, sowohl des Juden zuerst als auch des Griechen;
	Herrlichkeit aber und Ehre und Frieden jedem, der das Gute wirkt, sowohl dem Juden zuerst als auch dem Griechen.
	\attrib{\bibverse{Romans}{2}{5-10}, \Cite{elb}}
\end{quote}

Wir können also zusammenfassend sagen:
Beim letzten Gericht wird jeder Mensch nach seinen Taten gerichtet. Die Gerechten werden leben und die, die Gott ablehnen werden vernichtet.
Über diese Tatsache gibt es in der jüdischen Literatur bis ins erste Jahrhundert keine Diskussion, nicht einmal im Brief des Paulus an die Römer.
Wie passt das Gericht nach Werken und die `Gerechtfertigung' durch Glauben zusammen?\footnote{
	Innerhalb von 2 Kapiteln in Römer 2-3!
}
Um diese Frage beantworten zu können wenden wir uns jetzt den griechischen Wörtern von der Wurzel δικαι zu: δικαιος, δικαιοσυνη, und δικαιοω.
Wir werden sehen, dass sich der scheinbare Widerspruch nicht auflöst, indem wir das Gericht nach Werken umdefinieren, sondern indem
wir die δικαι-Wörter, und damit das Konzept der `Gerechtfertigung' im Neuen Testament besser verstehen.

\section{Die δικαι-Wörter}
Die δικαι-Wörter, insbesondere in den paulinischen Briefen, sind in der aktuellen wissenschaftlichen Literatur äußerst
umstritten. Bei der Frage `was meint Paulus mit Gerechtfertigung' existieren mindestens die folgenden Ansicht:
`Gerechtsprechung'\footnote{
	Die klassische protestantische Interpretation. Siehe zum Beispiel \Cite{both-judge-and-justifier}, \Cite[82ff]{nicnt-romans}.
}, `Deklaration als Mitglied des Bundes'\footnote{
	Verbunden mit der New Perspective on Paul, insbesondere gehalten von NT Wright.
}, `Wiederherstellung der Bundesbeziehung, verbunden mit der Hoffnung auf Freispruch'\footnote{
	\Cite{inhabiting-the-cruciform-god}. Diese Sicht ist mit NT Wrights Sicht verbunden - aber Wright schließt einen transformativen
	Teil in der Bedeutung von δικαιοω explizit aus (\Cite{paul-and-the-faithfulness-of-god}), während Gorman ihn mit einbezieht.
}.
Anderswo ist `justification inseparable from glorification'\footnote{
	\Cite[197]{paul-a-new-covenant-jew}. Pitre, Barber, und Kincaid sind explizit der Meinung, dass Gerechtfertigung ausschließlich als legale
	Transaktion nicht ausreichend ist.
}.

Bei der Betrachtung jedes Wortes ist es fundamental wichtig, zwischen der invarianten Bedeutung - also der Bedeutung, die ein Wort in jeden Kontext einbringt - und
der Interpretation des Wortes in einem Kontext - also wie die invariante Bedeutung mit dem Kontext zusammen interpretiert wird - zu unterscheiden.
Natürlich reichen diese Seiten nicht aus, um mit allen Positionen zu den δικαι-Wörtern im Detail zu interagieren.
Stattdessen möchte ich zeigen, dass
(a) die δικαι Wortgruppe in ihrer invarianten Bedeutung nicht auf eines der Extreme Ontologie oder Forensik beschränkt ist, und
(b) dass eine genaue Untersuchung des Kontexts nötig ist, um zwischen verschiedenen Interpretationen zu unterscheiden.
Ich werde auch eine Hypothese dafür geben, wie diese Disambiguierung im allgemeinen vorgenommen werden kann.

Das linguistisch komplexeste der δικαι-Wörter ist die verbale Form. Im Gegensatz dazu sind das adjektiv δικαιος und das deadjektivische Substantiv δικαιοσυνη deutlich einfacher zu analysieren:
\begin{exe}
	\ex \label{ex:dikaios-open}
	{
		ὁ δὲ ἄνθρωπος, ὃς ἔσται δίκαιος, ὁ ποιῶν κρίμα καὶ δικαιοσύνην, [...], δίκαιος οὗτός ἐστιν \\
		Und der Mensch, der gerecht ist, der Recht und Gerechtigkeit ausübt, [es folgt eine Liste von bösen Dingen, die dieser Mensch nicht tut], gerecht ist er.
	}
	\attrib{\bibverse{Ezekiel}{18}{5-9} LXX}
\end{exe}
An diesem Beispiel sehen wir gut, wie das Adjektiv und das Substantiv parallel verwendet werden und beide einen Menschen charakterisieren, der richtig handelt.

\begin{exe}
	\ex \label{ex:dikaios-gradable}
		{Δίκαιος σὺ ὑπὲρ ἐμέ \\
		 `du warst gerechter als ich'
		}
	\attrib{\bibverse{ISamuel}{24}{18} = LXX 1. Buch der Könige 24:18}
\end{exe}
Hier wird sichtbar, dass δικαιος für Vergleiche verwendet werden kann.
Adjektive, die für Vergleiche verwendet werden können nennt man auch \textit{gradierbar}: sie können zu bestimmten Graden zutreffen.
Beispiele im deutschen sind `groß' und `feucht', aber zum Beispiel nicht `tot'.
Gradierbare Adjektive haben einige besondere linguistische Eigenschaften, insbesondere hat die Art der \textit{Skala}, die sie beschreiben einen großen Einfluss auf ihr
syntaktisches Verhalten und ihre Interpretation.\footnote{
	Für eine genauere Einführung, siehe \Cite{scale-structure}.
}
Vergleichen wir beispielsweise diese deutschen Beispiele:
\begin{exe}
\ex\label{ex:open-closed} 
\begin{xlist}
	\ex[]{Max ist nicht groß, aber größer als Moritz.}
	\ex[]{\label{ex:feucht-gradable}Nach dem Schwimmen ist meine Badehose feuchter als mein Handtuch.}
	\ex[\texttt{\#}]{\label{ex:upper-half-open}Meine Hand ist nicht feucht, aber es ist etwas Wasser auf ihr.\footnote{\texttt{\#} bedeutet: Dieses Beispiel ist grammatisch erlaubt, aber in seiner Bedeutung widersprüchlich.}}
	\ex[*] {\label{ex:non-gradable}Bob ist nicht lebendig, aber Alice ist töter.\footnote{* bedeutet: Dieses Beispiel ist ungrammatisch.}}
\end{xlist}
\end{exe}
Wir sehen zunächst, dass manche adjektive nicht gesteigert und verglichen werden können: `tot' ist nicht gradierbar, und Beispiel \ref{ex:non-gradable} ist ungrammatisch.
Andere adjektive sind steigerbar und auf einer \textit{offenen Skala}. Das eine Person `groß' ist, entscheidet sich nur im Vergleich mit anderen Personen:
Sind Max und Moritz im selben Raum, so ist in diesem Kontext Max groß, und Moritz klein. Es gibt keine Untergrenze und Obergrenze für das `Großsein'.
Man sagt, dass die Skala beidseitig offen ist.
Beispiel \ref{ex:upper-half-open} zeigt, dass sich das deutsche adjektiv `feucht' anders verhält.
Dieses adjektiv ist gradierbar, wie \ref{ex:feucht-gradable} zeigt, aber es gibt eine fixe Untergrenze für `Feuchtheit'.
Sobald auch nur etwas Wasser auf einer Oberfläche ist, ist diese Oberfläche automatisch feucht.

\begin{exe}
\ex\label{ex:ambiguous-scale} 
\begin{xlist}
	\ex[]{\label{ex:dry-ambiguous}Meine Hände sind trocken.}
	\ex[]{\label{ex:dry-closed} Ich habe meine Hände gerade abgetrocknet und sie sind trocken.}
	\ex[]{\label{ex:dry-open} Meistens sind meine Hände trocken, aber unter Stress schwitzen sie schnell.}
\end{xlist}
\end{exe}

Beispiel \ref{ex:ambiguous-scale} schließlich zeigt, dass ein und dasselbe Wort `trocken' je nach Kontext verschiedene Skalenstrukturen realisieren kann:
In \ref{ex:dry-closed} ist die Skala geschlossen: eine individuelle Situation ist im Blick, und `meine Hände sind trocken' bedeutet
`es befindet sich aktuell keine Flüssigkeit auf ihnen' - das Maximum an Trockenheit ist in diesem Moment tatsächlich erreicht.
In \ref{ex:dry-open} wird eine allgemeine Qualität der Hände beschrieben.
Verglichen mit den Händen anderer Menschen sind meine Hände üblicherweise trocken.
Trotzdem können meine Hände in speziellen Situationen naß sein. Die Skala ist also offen, es wird nicht über ein Maximum an Trockenheit geredet.
Zuletzt ist die Aussage in \ref{ex:dry-ambiguous} unklar. Das deutsche Wort `trocken' hat weder eine offene Skala (allgemeine Qualität), noch eine geschlossene Skala
(Zustand in einer bestimmten Situation) als Teil seiner invarianten Bedeutung.

Mit dieser Theorie kommen wir nun zurück zu den δικαι-Wörtern.
In \ref{ex:dikaios-gradable} sahen wir bereits, dass δικαιος gradierbar ist.
Die gesamte Wortgruppe markiert ihre Skalenstruktur nicht, analog zum deutschen Wort `trocken'.\footnote{
	Hier liegt eines der Probleme, in deutscher oder englischer Sprache über die `Gerechtfertigung' zu reden:
	Im Deutschen und Englischen haben wir kein Verb, das in seiner Skalenstruktur so flexibel ist wie das griechische δικαιοω.
	Wir müssen also beim Übersetzen eine Entscheidung darüber treffen, ob das griechische Wort geschlossen oder offen zu interpretieren ist.
}
Das Adjektiv δικαιος kann eine offene Skala haben, und die grundsätzliche Art eines Menschen beschreiben, wie in \ref{ex:dikaios-open} oder in folgendem Beispiel:
\begin{exe}
	\ex
		{
			ἦσαν δὲ δίκαιοι ἀμφότεροι ἐναντίον τοῦ θεοῦ, πορευόμενοι ἐν πάσαις ταῖς ἐντολαῖς καὶ δικαιώμασιν τοῦ κυρίου ἄμεμπτοι. \\
			Beide aber waren gerecht vor Gott und wandelten untadelig in allen Geboten und Satzungen des Herrn.
		}
		\attrib{\bibverse{Luke}{1}{6}\footnote{
			Es ist wahr, dass die `Gebote und Satzungen des Herrn' einen Standard bilden, der die Skala nach oben abschließen könnte.
			In dieser Stelle im Diskurs wird allerdings der Charakter von Zacharias und Elisabeth beschrieben,
			keine individuelle Situation, in der ihr Handeln in Frage gestellt wird.
			Stattdessen ist ihr `Wandeln in den Geboten' ebenso eine allgemeine Charaktereigenschaft - aus theologischen Gründen wollen
			wir diese Verse so oder so nicht verstehen als wären Zacharias und Elisabeth sündlos - also in jeder individuellen Sitution ihres Lebens gerecht.
	}}
\end{exe}
Ein Beispiel, in dem die Skala geschlossen ist findet sich in \bibverse{IKings}{8}{31-32} (=LXX 3. Buch der Könige 8:31-32),
wo eine individuelle Sünde von einem Menschen gegen einen anderen in Sicht ist, und der Unschuldige δικαιος genannt wird.
Es ist nicht der Charakter der Personen, sondern eine individuelle Tat in Frage.\footnote{
	δικαιος wird sehr selten in dieser geschlossenen Form verwendet, viel häufiger taucht es mit einer offenen Skala auf.
	Der Grund für diese Tatsache ist wahrscheinlich, dass αθωος (unschuldig) zur Verfügung steht, das für eine geschlossene Skala markiert ist.
	Wir erwarten, und sehen in den Daten, dass das unmarkierte δικαιος nur selten in einer Art und Weise verwendet wird, die vom markierten αθωος besser gefüllt wird.
	Ein gutes minimales Beispiel bildet \bibverse{Matthew}{27}{24}.
	Pilatus ist unschuldig an einer individuellen Tat (geschlossene Skala), sein allgemeiner Charakter (offene Skala) ist hier nicht im Spiel und αθωος wird verwendet.
	Die verbale Form αθωοω ist deutlich seltener zu finden (nicht im NT, nur LXX), und damit korreliert wahrscheinlich die Tatsache, dass δικαιοω häufiger für die
	geschlossene Skala verwendet wird als sein Adjektiv. Siehe für ein klares Beispiel \ref{ex:dikaiow-closed}.
}
Das Substantiv δικαιοσυνη unterscheidet sich kaum vom Adjektiv.
Wenn wir nun zur verbalen Form δικαιοω kommen ensteht ein ähnliches Bild, aber die offene Skala wird häufiger verwendet als beim Adjektiv:
\begin{exe}
	\ex\label{ex:dikaiow-closed}
	\begin{xlist}
		\ex \label{ex:dikaiow-not-transformational} {
			Ἐὰν δὲ γένηται ἀντιλογία ἀνὰ μέσον ἀνθρώπων [...] καὶ δικαιώσωσιν τὸν δίκαιον καὶ καταγνῶσιν τοῦ ἀσεβοῦς \\
			Wenn es zu einem Rechtsstreit zwischen Männern kommt [...], und sie sprechen den Gerechten gerecht, und den Schuldigen für unrecht befinden
		}
		\attrib{\bibverse{Deuteronomy}{25}{1} LXX}
		\ex \label{ex:dikaiow-forensic-real} {
			ἐκ γὰρ τῶν λόγων σου δικαιωθήσῃ καὶ ἐκ τῶν λόγων σου καταδικασθήσῃ \\
			denn aus deinen Worten wirst du gerecht befunden werden, und aus deinen Worten wirst du verdammt werden
		}
		\attrib{\bibverse{Matthew}{12}{37}\footnote{
				Es ist hier (a) das Gericht im Blick, (b) stehen individuelle Worte im Fokus, anhand derer die Menschen gerichtet werden.
		}}
	\end{xlist}
	\ex\label{ex:dikaiow-open}
	\begin{xlist}
		\ex {
				Ἄρα ματαίως ἐδικαίωσα τὴν καρδίαν μου καὶ ἐνιψάμην ἐν ἀθῴοις τὰς χεῖράς μου \\
				Jetzt habe ich umsonst mein Herz gerecht gemacht, und in Unschuld meine Hände gewaschen!
		}
		\attrib{\bibverse{Psalms}{73}{13} = LXX Psalms 72:13}
		\ex \label{ex:dikaiow-transformational} {
			καὶ ταῦτά τινες ἦτε ἀλλʼ ἀπελούσασθε, ἀλλʼ ἡγιάσθητε, ἀλλʼ ἐδικαιώθητε \\
			Und das [Sünder] sind einige von euch gewesen, aber ihr seid abgewaschen, aber ihr seid heilig gemacht, aber ihr seid gerecht gemacht
		}
		\attrib{\bibverse{ICorinthians}{6}{11}}
	\end{xlist}
\end{exe}
In den Beispielen \ref{ex:dikaiow-closed} ist die Skala geschlossen, und das Verb muss dementsprechend interpretiert werden:
\ref{ex:dikaiow-not-transformational} erlaubt keine Interpretation außer,
dass die Person in Frage in diesem individuellen Rechtstreit Recht erhält oder für gerecht befunden wird.
In beiden Beispielen in \ref{ex:dikaiow-open} muss der allgemeine Charakter oder der generelle Zustand der Person (bzw. des Herzens) in Frage sein,
da keine individuelle Tat oder ein spezieller Streit das Thema ist.
Gerade in \ref{ex:dikaiow-transformational} ist keine Interpretation möglich außer, dass δικαιοω hier parallel zu `abgewaschen' und `heilig gemacht' steht:
für eine Transformation des allgemeinen Charakters der Person.
Die Interpretationen in \ref{ex:dikaios-open} sind notwendigerweise inchoativ (sie markieren den Beginn eines Zustands).
Die Interpretationen mit geschlossener Skala in \ref{ex:dikaiow-closed} dagegen beschreiben das ausdrücken oder feststellen einer
forensischen Realität durch das Aufrichten oder für gerecht Erklären des Gerechten.\footnote{
	Vergleiche auch mit \Cite{hebrew-verbs-of-judgement}.
	Die Situation ist bei den hebräischen und griechischen Wörtern aus dieser Domäne sehr ähnlich.
}
Zusammenfassend können wir also sagen, dass die δικαι-Wörter unmarkiert für die Skalenstruktur sind.
Der Kontext muss disambiguieren, ob
(a) eine geschlossene Skala, also eine individuelle Tat, und damit die Schuld oder Unschuld oder die Erfüllung eines fixen Rechtsstandards, oder
(b) eine offene Skala, also der Charakter oder die allgemeine Disposition einer Person beschrieben wird.

Um es theologischer auszudrücken:
Eine Theorie, die eine fixe theologische Bedeutung für die δικαι-Wörter fixieren möchte ist zum Scheitern verurteilt.
In jedem Fall ist der Kontext der individuellen Stelle nötig, um zwischen grundlegend verschiedenen Interpretationen unterscheiden zu können.
Werden den Geretteten ihre Sünden nicht angerechnet? Ja.
Formen die Geretteten im Christus den neuen Bund, dessen Mitglieder in rechter Beziehung zu Gott stehen? Ja.
Werden die Gläubigen aus ihrem Leben der Sünde heraus gerissen, und vom Geist befähigt Gott wohlgefällig zu leben? Ja.
Aber keine dieser Aussagen darf mit dem Wort `Gerechtfertigung' gleichgesetzt werden. Jeder Text, und selbst jede Phrase in einem Text muss für sich stehen und gelesen werden
um zu einer belastbaren Theologie zu gelangen.

\section{Gerechtigkeit und das Kreuz im Brief an die Römer}
In seinem Brief an die Römer beschreibt der Apostel Paulus so ausführlich wie kaum anderswo wie Gott die Menschen
durch das Wirken seines Sohne aus ihrer miseren Lage holt.
Da wir hier in einem anhaltenden Argument über das letzte Gericht, Jesu Tod am Kreuz, unseren neuen Frieden mit Gott, und unser neues Leben in Christus hören
möchte ich im folgenden zeigen, wie das neue Verständnis der δικαι-Wörter hilft, den Römerbrief konsistent zu interpretieren,
und wie in diesem Verständnis Jesu Tod in stellvertretender Übernahme unserer Schuld keine Rolle spielt.

\subsection{Römer 2}
Mit \bibverse{Romans}{2}{} beginnen wir sofort in einer äußerst umstrittenen, aber gleichzeitig oft ignorierten Passage des Briefe.
In Abschnitt \ref{sec:gericht-nach-werken} sahen wir bereits, dass dieses Kapitel eine der Stellen im Neuen Testament ist, in der das endgültige Gericht
nach Werken beschrieben wird.
Die größte Interpretative Frage zu diesem Kapitel ist nicht, was es bedeutet, sondern ob es konstruiert ist und damit eine hypothetische Situation beschreibt:
ein Gericht nach Werken, vor dem die Gläubigen gerettet werden.
Ist es also  Paulus' Argument, dass vor diesem Gericht, wenn es wirklich nach Werken geschieht, kein Mensch gerecht ist?

Diese Sicht ist äußerst unplausibel. Sicherlich ist die Klasse der gerechten Menschen in Texten über das letzte Gericht
klein: die meisten Menschen, selbst die meisten Juden, sind ungerecht und werden im letzten Gericht vernichtet.
Doch die klassische Position, dass kein Mensch vor Gottes Gericht gerecht ist findet sich nirgends sonst\footnote{
	\bibverse{Luke}{1}{6} war bereits zitiert: hier sind zwei Menschen `gerecht und wandeln tadellos in allen Gesetzen und Satzungen des Herrn'.
	Hat sich Lukas versprochen?
}, und erzeugt auch keine konsistente Interpretation dieses Kapitels.
Um das zu zeigen, betrachten wir zunächst v13:
\begin{quote}
	es sind nämlich nicht die Hörer des Gesetzes gerecht vor Gott, sondern die Täter des Gesetztes werden gerecht gesprochen werden.
	\attrib{\bibverse{Romans}{2}{13}}
\end{quote}
Es ist in der ganzen Passage (a) das letzte Gericht, und (b) ein klarer Standard in Sicht: Gottes Wille im Gesetz offenbart.
Damit ist die Skalenstruktur des Verbs δικαιοω geschlossen und eine entsprechende Interpretation wie hier angegeben ist nötig.
Das Ausführen des Gesetzes (der Torah) ist nötig, um von Gott im Gericht als gerecht befunden zu werden.
\begin{quote}
	Denn wenn Nationen, die von Natur aus kein Gesetz haben\footnote{
		Diese Übersetzung ist mit großer Sicherheit korrekt, kontra \Cite{nicnt-romans}.
		Die griechische Phrase ist ἔθνη τὰ μὴ νόμον ἔχοντα φύσει τὰ τοῦ νόμου ποιῶσιν, und die Frage ist, ob φύσει die erste Phrase `Nationen, die das Gesetz nicht haben',
		oder die zweite Phrase `sie tun das Gesetz' modifiziert.
		(a) In adverbialer Verwendung modifiziert φύσει Zustände, keine Ereignisse (in NT und LXX: \bibverse{Galatians}{2}{15} mit Null-Copula, \bibverse{Galatians}{4}{8}, \bibverse{Ephesians}{2}{3}, \bibverse{Wisdom}{13}{1} mit Null-Copula).
		(b) Die Idee, dass Heiden von Natur aus dem Gesetz entsprechend handeln ist inkonsistent mit \bibverse{Galatians}{2}{15}. Heiden sind von Natur aus Sünder, und haben von Natur aus das Gesetz nicht.
	}, das Gesetz tun,
	so sind diese, die kein Gesetz haben, sich selbst ein Gesetz.
	Sie beweisen, dass das Werk des Gesetzes in ihren Herzen geschrieben ist
	\attrib{\bibverse{Romans}{2}{14-15}}
\end{quote}
Es gibt also Heiden, die zwar von Natur aus kein Gesetz haben, aber doch die Dinge des Gesetzes tun.
So beweisen sie, dass Gottes Gesetz auf ihren Herzen geschrieben ist - die prophetische Hoffnung, die wir schon in Kapitel \ref{chap:prophetische-hoffnung} gesehen haben
und die jetzt durch den Neuen Bund im Christus Jesus offenbart ist.
Und aus diesem Grund können diese Heiden die Juden mit denen Paulus sein Argument führt anklagen: Sie tun das Gesetzt, in Christus, und werden deshalb im Gericht bestehen.
Die Juden ruhen sich darauf aus, die Torah zu kennen aber erlangen doch Abseits von Christus keine Gerechtigkeit.
Diese Verse zeigen, dass es für Paulus sehr wohl möglich ist für einen Menschen, das Gesetz zu erfüllen - nämlich in Christus.
Genau zu diesem Punkt werden wir in der Diskussion über \bibverse{Romans}{8}{3} zurückkehren: ``die Rechtsforderung des Gesetzes wird erfüllt in uns, die wir nicht nach dem Fleisch, sondern nach dem Geist wandeln''.

Ein sehr ähnliches Argument finden wir am Ende des Kapitels:
Paulus beschuldigt Juden, die sich auf das Gesetz stützen und doch Böses tun (v17-24), und setzt im Kontrast dazu den Heidenchristen, beschnitten im Herz statt am Körper,
dem diese Beschneidung tatsächlich etwas nützt:
\begin{quote}
	Wenn nun der Unbeschnittene die Rechtsforderungen des Gesetzes befolgt, wird nicht sein Unbeschnittensein für Beschneidung gerechnet werden,
	und das Unbeschnittensein von Natur, das das Gesetz erfüllt, dich richten, der du mit Buchstaben und Beschneidung ein Gesetzesübertreter bist?
	\attrib{\bibverse{Romans}{2}{26-27}, \Cite{elb}}
\end{quote}
Diese Argument ist absolut sinnlos, wenn Paulus hier über eine hypothetische Person redet, in dem Wissen dass sie nicht existiert, und
dass er im nächsten Kapitel erklären wird, dass diese Person nicht existiert.
Paulus legt hier eine sehr reale Situation auf den Tisch: der Heide, der abseits der körperlichen Beschneidung die Forderung des Gesetzes erfüllt
klagt den untreuen Juden an, den Paulus zur Umkehr bewegen will.\footnote{
	Moos Argument `But Paul would not describe a Christian as having been granted membership among the saved as a result of obeying the law' \Cite{nicnt-romans}
	ist zirkulär, und verliert seine Kraft wenn wir im folgenden sehen werden, dass `Gerechtfertigung aus Glauben' nicht `im letzten Gericht unabhängig von Werken
	gerecht gesprochen' bedeutet.
	Wenn Paulus das letzte Gericht so versteht wie jeder andere jüdische Text seiner Zeit, und so wie seine eigene Beschreibung auf den ersten Blick aussieht,
	dann ist `Christen werden im letzten Gericht gerecht erklärt, weil sie wirklich gerecht waren' exakt Paulus' Meinung.
}
Die Tuer des Gesetzes werden in Gottes Gericht bestehen.
Es stellen sich jetzt zwei Fragen: erstens, wie kann es sein, dass die Juden die das Gesetz haben nicht dadurch Tuer des Gesetzes sind.
Zweitens, wie kann es sein, dass die Heiden ohne Gesetz Tuer des Gesetzes werden?\footnote{
	Für eine detailliertere Diskussion über verschiedene Meinungen zu Römer 2 und eine anhaltendere Verteidigung der realen Position
	siehe \Cite[215ff]{judgement-and-justification-in-early-judaism-and-the-apostle-paul}.
}

\subsection{Römer 3:9-20}
Die tatsächliche Grundlage dafür, dass viele Kommentatoren Kapitel 2 auf die eine oder andere Art und Weise als hypothetisch abschreiben ist,
dass sie einen Widerspruch zwischen Paulus' Aussagen in Kapitel 3 sehen. Dieser Stelle wenden wir uns nun zu.
\begin{quote}
	Was nun? Haben wir einen Vorzug? Durchaus nicht! Denn wir haben sowohl Juden als auch Griechen vorher beschuldigt, alle unter der Sünde zu sein,
	wie geschrieben steht: ``Da ist kein Gerechter, auch nicht einer; da ist keiner, der verständig ist; da ist keiner, der Gott sucht.
	Alle sind abgewichen, sie sind allesamt untauglich geworden; da ist keiner, der Gutes tut, da ist auch nicht einer.''
	[...]
	% ``Ihr Schlund ist ein offenes Grab; mit ihren Zungen handelten sie trügerisch.'' ``Viperngift ist unter ihren Lippen.''
	% ``Ihr Mund ist voll Fluchens und Bitterkeit.''
	% ``Ihre Füße sind schnell, Blut zu vergießen; Verwüstung und Elend ist auf ihren Wegen, und den Weg des Friedens haben sie nicht erkannt.''
	% ``Es ist keine Furcht Gottes vor ihren Augen.''
	Wir wissen aber, dass alles, was das Gesetz sagt, es denen sagt, die unter dem Gesetz sind,
	damit jeder Mund verstopft wird und die ganze Welt dem Gericht Gottes verfallen ist.
	Darum: Aus Gesetzeswerken wird kein Fleisch vor ihm gerechtfertigt werden;
	denn durchs Gesetz \textlangle kommt\textrangle\ Erkenntnis der Sünde.
	\attrib{\bibverse{Romans}{3}{9-20}, \Cite{elb}}
\end{quote}
Im Kontrast zu \bibverse{Romans}{2}{13} ``die Täter des Gesetzes werden gerecht gesprochen werden'' bildet v20 eine Interpretatives Rätsel.
Besonders die protestantische Position löst dieses Problem häufig auf, indem 2:13 hypothetisch gelesen wird: Im Prinzip würden die Täter des Gesetzes schon
gerecht gesprochen werden, aber Paulus zeigt mit seinen Zitaten aus dem Alten Testament, dass kein einziger Mensch diesen Standard tatsächlich erfüllt.
Da ich im letzten Abschnitt dafür plädiert habe, dass Römer 2 nicht hypothetisch uminterpretiert werden kann muss sich dieses Paradox also in 3:20 auflösen.
Die primäre Erklärung, mit der sich dieses Paradox auflöst ist die Feststellung, dass δικαιοω in v20 auf einer offenen Skala agiert, dass Paulus
also über den allgemeinen Zustand der Menschen, nicht die individuellen Sünden des Lebens am letzten Gericht spricht.
Entsprechend ist meine Interpretation, dass `aus Werken des Gesetzes kein Mensch gerecht wird'.
Für diese Sicht sprechen folgende Punkte:

Erstens: Die Bibelstellen, die Paulus für sein Argument zu Rate zieht sagen nicht aus, dass keine einzige Person vor Gottes Gericht bestehen kann.
Sie zeigen den natürlichen Hang aller Menschen - auch der Juden - zur Sünde.
\bibverse{Ecclesiastes}{7}{21} sagt aus, dass kein Mensch so gerecht ist, dass er nie sündigt.
Aber es ist in der jüdischen Literatur klar, dass `Gericht nach Werken' nicht `jeder der auch nur einmal etwas Böses getan hat wird untergehen' bedeutet.
Sündlose Perfektion ist nicht der Rahmen in dem Gericht nach Werken zu verstehen ist.\footnote{
	Siehe zum Beispiel \Cite[CD-A Col II = 4Q266 2 II]{tdss-se}: ``God loves knowledge; he has established wisdom and counsel
	before him; prudence and knowledge are at his service; patience is his and
	abundance of pardon, to atone for those who repent from sin; however,
	strength and power and a great anger with flames of fire by the \textlangle hand\textrangle\ of all
	the angels of destruction against those turning aside from the path and abominating the precept''
	Selbst in der sehr streng legalistischen Gemeinschaft in Qumran ist axiomatisch akzeptiert, dass Menschen umkehren und sich der Wahrheit zuwenden können.
}
Einige der anderen zitierten Stellen beschreiben, dass Menschen vor Gott gerecht oder unschuldig ist, und Menschen deswegen Gnade brauchen.\footnote{
	\bibverse{Psalms}{143}{2} (der Author), \bibverse{Isaiah}{59}{7-8} (Israel).
	Diese Aussage widerspricht meiner Interpretation nicht: selbstverständlich sündigt jeder Mensch, und natürlich ist es Gottes Gnade,
	die dem Gerechten seine Schuld vergibt.
	Mein Punkt ist: wir sollten solch negative Anthropologie nicht soweit strapazieren, dass jede Gerichtsszene
	zu einer Farce wird, in der niemand auf der Seite Gottes steht.
	Sieht Jesus in \bibverse{Matthew}{25}{31-46} oder \bibverse{John}{5}{28-29} Szenen, in denen er von zwei Seiten im Gericht spricht -
	aber eine der Seiten existiert überhaupt nicht, weil niemand gerecht ist?
	Ist es nicht wahrscheinlicher, dass in poetischem Material der Sprecher eine erhöhte Sicht seiner Unwürdigkeit vor Gott wiedergibt
	und Gott aus Gnade einen Maßstab an Menschen setzt, der tatsächlich realistisch erfüllbar ist?
	Die Tatsache dass das schwarz-weiße Aufteilen der Menschen in Gerechte und Ungerechte im letzten Gericht so häufig formuliert
	wird macht dieses Verständnis deutlich einfacher.
}
In den meisten zitierten Psalmen schreit der Sprecher zu Gott einzugreifen, indem er die ungerechten und gottlosen Unterdrücker bestraft.
Gleichzeitig ist explizit, dass der Sprecher sich selbst auf der Gerechten Seite der Medaillie sieht.
In allen diesen Texten ist klar, dass nicht jeder einzelne Mensch ungerecht im Sinne des Psalmes ist.\footnote{
	Vergleiche \bibverse{Psalms}{14}{1-3} mit v5; \bibverse{Psalms}{5}{10} mit v13; \bibverse{Psalms}{140}{4} mit v14; \bibverse{Psalms}{36}{8} mit v11.
	In \bibverse{Psalms}{10}{} spricht der `Elende' und bittet Gott um Rettung vor dem Gottlosen.
	Die Implikation ist, dass der Elende im Sinne des Psalmes
	nicht ebenso Gottlos ist.
	In \bibverse{Proverbs}{1}{16} nutzt der Lehrer die Aussage, um seinen Schüler von den Ungerechten fern zu halten. Die Aussage ist wiederum nicht,
	dass jeder Mensch in der selben Ungerechten Kategorie ist.
}

Zweitens: Paulus argumentiert nicht, dass `alle Sünden getan haben', sondern dass `alle unter der Sünde sind' - selbst die Juden.
Genau das zeigen die Zitate aus dem Alten Testament wirklich: jeder Mensch, Heide oder Jude, lebt als Mensch unter der Sünde.
Jeder einzelne Mensch sündigt tatsächlich, und niemand kann Gottes Perfektion entsprechen.
Genau das schreibt auch Paulus selbst: ``Denn wir haben sowohl Juden als auch Griechen vorher beschuldigt, alle unter der Sünde zu sein'' (3:9).
Paulus will nicht beweisen, dass `alle Menschen aus Sünde an Gott schuldig geworden sind', oder dass `jeder einzelne Mensch vor Gottes Gericht ungerecht ist'.
Das Problem, unter das Paulus hier Juden ebenso wie die Griechen einschließt ist, unter der Kraft der Sünde versklavt zu sein.\footnote{
	Ein Punkt, den Paulus immer wieder verwendet, beispielsweise in \bibverse{Romans}{6}{17}.
}

Drittens: Es ist in Kapitel 2 bereits klar, dass das Herz einer Person zu allerlei Sünden führt:
``Nach deiner Störrigkeit und deinem unbußfertigen Herzen aber häufst du dir selbst Zorn auf für den Tag des Zorns''.
Das ursprüngliche Problem ist das Herz, die Störrigkeit des Menschen, sein Wesen selbst -
nicht bloß die Tatsache, dass Menschen Schuldig sind und es eine Lösung braucht, damit Gott Sündern die Umkehren wollen vergeben kann.
Es ist die Natur des Menschen an sich in Frage.

Viertens: Es ist auch hilfreich zu verstehen, wie Paulus anderswo mit diesem Problem umgeht:
``denn alle haben gesündigt und es fehlt ihnen die Herrlichkeit Gottes''(3:23) zeigt deutlich,
dass das Problem größer ist als Schuld. Etwas wirkliches, Anfassbares fehlt dem Menschen in seinem sündigen Zustand.

Fünftens: Die Übersetzung `dem Gericht Gottes verfallen' in v19 ist nicht notwendig.
Das griechische ὑπόδικος τῷ θεῷ bezieht sich nicht explizit auf Gottes Gericht, selbst wenn die Elberfelder so übersetzt.\footnote{
	Und klassische Interpretationen wie \Cite{nicnt-romans} lesen hier natürlich eine Szene im Gerichtssaal.
}
Alle Menschen sind `haftbar' vor Gott, allen Menschen - selbst den Juden - ist der Mund gestopft, denn alle Menschen sind in Sünde.
Ob diese Tatsache im Kontext des letzten Gerichts zu verstehen ist oder nicht wird durch diesen Vers nicht markiert.

Zusammenfassend ergeben diese Punkte eine Szene, in der der allgemeine Zustand der Menschheit - der Heiden und Juden - in Frage steht.
Alle Menschen haben ein Herz-Problem\footnote{
	Diese Bezeichnung geht auf \Cite[179]{paul-a-new-covenant-jew} zurück.
	Ich gehe zwar in meiner linguistischen Analyse der δικαι-Wörter anders vor, stimme aber ihren theologischen Argumenten im Großen und Ganzen zu.
}, und dieses Problem muss gelöst werden.
Da keine individuelle Situation im Kontext steht, und Paulus eindeutig über den allgemeinen Zustand aller Menschen unter der Sünde redet
bleibt die Skala für δικαιοω in 3:20 offen und die Gerechtigkeit die das Gesetz nicht erzeugen kann ist die (echte, ontologische) Gerechtigkeit des Herzens.
``Menschen werden nicht aus den Werken des Gesetzes vor Gott gerecht gemacht, \textit{denn durch das Gesetz ist Erkenntnis der Sünde}.'' (3:20).
Das Gesetz erfüllt also nicht die Funktion, Menschen gerecht zu machen, denn stattdessen führt es nur zur Erkenntnis der Sünde.
Das Gesetz ist ein Hilfsmittel, ein Maßstab, an dem Gutes und Böses handeln gemessen werden kann.
Wer diese Gesetze tut wird leben (2:13), aber alle Menschen sind der Sünde verfallen (3:9-18),
und das Gesetz selbst kann den Menschen nicht dazu befähigen so zu handeln (3:20).
Paulus Aussage ist, dass das Gesetz kraftlos ist, Menschen gerecht zu machen.
Es zeigt Sünde auf, kann aber das eigentliche Problem nicht lösen.

Die Verse 2:1 - 3:20 bilden gewissermaßen ein Crux Interpretum und ein Gitter durch das hindurch der Rest des Römerbriefes gelesen wird.
Im Folgenden sehen wir, wie sich unsere Entscheidung in diesen Stellen im Rest des Briefes niederschlägt:
Wir haben ein klares Problem - Menschen sind unter Sünde, und auch das Gesetz ist nicht in der Lage, sie aus dieser Lage heraus gerecht zu machen.
Wenn unsere Interpretation hier korrekt ist sollte dieses Zentrale Problem immer wieder der Kern sein, um den sich Paulus weitere Argumente und seine Lösung dreht.

\subsection{Römer 3:21-26}
\begin{quote}
	Nun aber wurde abseits der Torah Gottes Gerechtigkeit offenbart, bezeugt von der Torah und den Propheten.
	Die Gerechtigkeit Gottes durch den Jesus-Christus-Glauben für alle die glauben.
	Es ist kein Unterschied: Denn alle haben gesündigt, und es fehlt ihnen die Herrlichkeit Gottes.
	Als Geschenk werden sie gerecht gemacht durch seine Gnade, durch den Loskauf in Christus Jesus.
	Ihn stellte Gott in seinem Blut öffentlich hin als Ort der Offenbarung durch Glauben
%	\footnote{
%		Wie genau diese Phrase zu verstehen ist, wird nicht direkt offensichtlich: ist (a) ιλαστηριον durch Glaube zugänglich,
%		oder ist (b) Jesu Treue zu Gottes Plan die Art und Weise, in der diese Offenbarung dargestellt wird?
%		Ich favorisiere die zweite Variante, aber da diese Phrase für unsere theologische Frage nicht relevant ist
%		habe ich hier eine sich nicht festlegende Übersetzung gewählt.
%	}
	zum Beweis seiner Gerechtigkeit im Blick auf das unbestraft lassen der früher geschehenen Sünden
	in der Nachsicht Gottes;
	zum Beweis seiner Gerechtigkeit in der jetzigen Zeit
	sodass er gerecht ist und aus dem Jesus-Glauben gerecht macht.
	\attrib{\bibverse{Romans}{3}{21-26}}
\end{quote}

Einige Kommentare zur Übersetzung an sich sind als aller erstes nötig.
\textbf{Torah} Das ``Gesetz'' ist im ganzen Kontext immer wieder das mosaische Gesetz. Dieses Torah ist es, die die Nationen von Natur aus nicht haben.
		Paulus Argument, gerade in Kapitel zwei, dreht sich immer wieder darum, dass das haben des Gesetzes nicht zum leben führt, sondern das tun.
		Ich sehe im Kontext des Römerbriefes keinen Grund, dieses Gesetz auf die ethnischen Grenzmarker zu beschrenken.

\textbf{Jesus-Christus-Glauben} Die Übersetzung der Phrase πίστεως Ἰησοῦ Χριστοῦ ist äußerst umstritten.
		Die klassischen Positionen sind (a) `der Glauben an Jesus Christus', die sogenannte objektive Sicht, oder (b) `die Treue Jesu Christus', die sogenannte Subjektive Sicht.
		Zu häufig wird die Diskussion über die Bedeutung einer griechischen Phrase als eine theologische Diskussion geführt.\footnote{
			Kontra \Cite[53]{porter-pitts} ``We cannot hope to solve such a delicate and sensitive issue as the meaning of the phrase πίστις Χριστοῦ simply by means of Greek linguistics,
			because there is more at stake – including exegesis and theology''.
			Aber wie kann unsere Exegese und Theologie begründet sein, wenn linguistische Fragen nicht linguistisch geklärt werden?
			Wenn Menschen sprechen und Authoren schreiben meinen sie etwas.
			Unsere Theologie in eine linguistische Frage zurück zu lesen ist eine sehr gefährliche Präzedenz.
		}
		Stattdessen ist es meine Position, dass die detaillierte linguistische Analyse primär gegenüber der Exegese oder gar der systematischen Theologie sein muss.
		Was ist also die beste Interpretation von πιστις Χρισοτου?
		Um diesen Artikel nicht um eine weitere komplexe Diskussion zu erweitern gebe ich hier nur die wichtigsten Argumente aus
		\Cite{grasso-pistis-christou} wieder.
		πιστις ist ein deverbales Substantiv. Solche Substantive können die Argumentstruktur des zugehörigen Verbes behalten.
		Die objektive Sicht bedeutet, dass das direkte Objekt des Verbes zum Genitivkonstrukt des Substantives wird.
		Wenn das Verb πιστευω allerdings ein direktes Objekt im Akkusativ annimmt, dann ist die Bedeutung immer `Glauben an eine Proposition'\footnote{
			\Cite[4.1.1]{grasso-pistis-christou}, cf. \bibverse{John}{11}{26}
		}.
		Da ein objektiver Genitiv diese Argumentstruktur beibehält, muss die Phrase dann `der Glauben, dass (eine unspezifizierte Präposition über) Christus wahr ist' bedeuten.
		Die (notwendige) Bedeutung `Jesus vertrauen' ist in dieser grammatischen Form unmöglich.
		Die subjektive Sicht ist ausgeschlossen, weil (a) das direkte Objekt (oder das indirekte Dativ-Objekt) direkt aus dem Kontext ersichtlich sein muss -
		aber das ist in den wichtigsten Stellen in denen die Phrase auftaucht nicht der Fall (b) es gibt keine Belege dafür, dass griechisch sprechende Kirchenväter
		die Phrase so verstanden haben.
		Die bessere Interpretation (der sogenannte `dritte Weg') ist der in der Übersetzung wiedergegebene Jesus-Christus-Glauben.
		Das ist das Glaubenssystem, dass Jesus Christus als seinen Inhalt hat.
		Eine parallele Formulierung hierzu finden wir in \bibverse{Philippians}{1}{2} συναθλοῦντες τῇ πίστει τοῦ εὐαγγελίου `zusammen für den Glauben des Evangeliums kämpfen'.
		Der Glaube hat das Evangelium als Inhalt - genauso hat der Jesus-Christus-Glauben Jesus Christus als Inhalt.

\textbf{gerecht machen} Es ist Sünde und fehlende Herrlichkeit im Blick, nicht `Übertretung in einer speziellen Situation' oder `Übertretung im Blick des letzten Gerichts'.
		Auch haben wir in den letzten Kapiteln bereits gesehen, wie immer wieder der Charakter der Person, nicht bloß die Sünde-als-einzelne-Tat, oder Sünde-als-Übertretung
		Paulus Anliegen ist.
		Die Skala bleibt also offen, und dementsprechend muss δικαιοω offen gelesen werden.

Der \textbf{Ort der Offenbarung} führt uns zu einer komplizierteren Diskussion. Ich vertrete hier insbesondere die Meinung von Porter in \Cite{mercy-seat-revelation}\footnote{
	Ähnlich wie Bailey in \Cite{jesus-as-the-mercy-seat}.
}.
Das griechische ἱλαστήριον wird in den folgenden Bedeutungen verwendet: (1) Der Deckel der Bundeslade (a) als Ort der Offenbarung\footnote{
	Das ist meine Interpretation. \bibverse{Exodus}{25}{17-22} (das ist die Bedeutung direkt bei der Einsetzung),
	\bibverse{Leviticus}{16}{2}, \bibverse{Numbers}{7}{89}.
	LXX \bibverse{Habakuk}{3}{2} ist eine klare Anspielung auf die Bundeslade (ἐν μέσῳ δύο ζῴων) und beschreibt durchweg eine Offenbarung Gottes an diesem Ort.
}
(b) als Gegenstand, der an Yom Kippur gereinigt wird\bibfoot{Leviticus}{16}{13-15},
(2) als allgemeiner Ort, an dem Opfer gebracht werden\footnote{
	\label{fn:hilstareion-in-iv-mac}
	\bibverse{Amos}{9}{1} muss sich auf die Obere Fläche eines Altares beziehen.
	\bibverse{Ezekiel}{43}{14-20}, wo es mehrere ἱλαστήρια gibt auf denen allerlei Opfer gebracht werden.
	Als letztes auch \bibverse{IVMaccabbees}{17}{22}.
	Diese Stelle wird oft als (die einzige) Stelle zitiert, in der ιλαστηριον die selbe Bedeutung wie in \bibverse{Romans}{3}{25}
	hat, nämlich metonymisch als Opfer, das die Sünden des Volkes wegnimmt.
	Siehe \Cite{mercy-seat-revelation} für eine detailliertere Auseinandersetzung:
	Der ιλαστηριον wird hier verwendet \textit{als Utensil}, mit dem die Brüder zu Tode gefoltert werden.
	Siehe auch \Cite[2,4]{ignatius-of-antioch-to-the-romans}, wo die wilden Tiere der Altar(θυσιαστήριον) sind, auf dem Ignatius geopfert wird (ἵνα διὰ τῶν ὀργάνων τούτων Θεοῦ θυσία εὑρεθῶ).
	Diese Stelle ist in der Formulierung sehr ähnlich zu \bibverse{IVMaccabbees}{17}{22}: διὰ [...] καὶ τοῦ ἱλαστηρίου τοῦ θανάτου αὐτῶν.
}
(3) als Weihgabe in der Heidnischen Welt.\footnote{
	Gemeint sind hier insbesondere Steelen, die als Versöhnungsgeschenk aufgestellt werden.
	Ein Beispiel ist \Cite[16.182]{antiquitates}.
}
Unsere Stelle im Römerbrief ist durchzogen von Wörtern der Offenbarung Gottes:
seine Gerechtigkeit ist offenbart, Jesus ist öffentlich hingestellt, zum Beiweis Gottes Gerechtigkeit.
Wir haben also direkt eine gute Korrelation zwischen Jesus-als-Offenbarung und dem Kontext.
Erzwungen wird diese Interpretation aber letztlich dadurch, dass alle anderen Vorschläge ausgeschlossen sind:
ιλαστηριον bedeutet einfach in keinem extenten Text `Sühneopfer', `Sühneort', `Sühnezeichen' oder ähnliches.
Es gibt \textit{nicht einen Text}, in dem diese Bedeutung vorkommt.\footnote{
	Wie in Fußnote \ref{fn:hilstareion-in-iv-mac} bereits angeschnitten ist die Situation in \bibverse{IVMaccabbees}{17}{22} ähnlich wie in \bibverse{Romans}{3}{25}.
	Es gibt eine bessere Interpretation, die zur Stelle passt und ihre Bedeutung nicht metonymisch (Deckel der Bundeslade, reinterpretiert als Sühneopfer) erzeugen muss.
}
Eine Übersetzung, in der der Deckel der Bundeslade metonymisch für ein Sühneopfer verwendet wird ist äußerst unplausibel,
einfach weil deutlich einfachere Interpretationen existieren.
Der Deckel der Bundeslade selbst ist gerade der Gegenstand, der in \bibverse{Leviticus}{16}{} gereinigt werden muss,
also ist auch eine Interpretation im Sinne von (1.b) nicht sinnvoll.
Die Bundeslade ist kein Altar im eigentlichen Sinne - sie ist der Wohnort Gottes, der von Unreinheit und Sünde befleckt wird.
Das Ziel des Dekontaminationstags ist es gerade, diesen Wohnort Gottes für seine Gegenwart zu säubern, wie wir in Abschnitt \ref{sec:yom-kippur} gesehen haben.
Was auch immer Paulus ausdrücken will, es ist mit Sicherheit nicht, dass Jesus der Ort ist der gereinigt werden muss.\footnote{
	Ein Argument, das selbst Porter nicht angibt, auch wenn wir für die selbe Position argumentieren.
}
Option (2) fällt aus einem ähnlichem Grund aus. Wenn Jesus der Altar ist, was wird dann geopfert? Jesus, also der Altar, selbst?
Option (3) ist in der aktuellen Wissenschaft recht unbeliebt, aber tatsächlich deutlich einfacher zu halten als (1a) oder (2).
Aber Paulus versucht hier ein größeres Problem zu lösen, als dass Menschen und Gott im Streit sind.
Was Jesus tut sollte ein Sünden-Problem lösen.\footnote{
	Für eine moderne Formulierung dieser Sicht (die ich am Ende doch ablehne), siehe \Cite[259ff]{lamb-of-the-free}.
}

Wenn wir ιλαστηριον in 3:25 als Deckel der Bundeslade \textit{als Ort der Offenbarung} (1) lesen, wie erklärt es dann die Passage im ganzen?
Zunächst sollten wir uns fragen, was JHWH im Alten Bund an diesem Ort offenbart hat.
Wenn Mose im Zelt der Begegnung mit Gott spricht, dann geschieht das an der Bundeslade\bibfoot{Numbers}{7}{89}.
Zu diesem Zeitpunkt ist der Großteil des Gesetzes an sich bereits offenbart - aber Mose erhält in seinem Gespräch mit Gott
\textit{Interpretationen und Anwendungen} des Gesetzes auf verschiedene Situationen.
In \bibverse{Numbers}{9}{} besteht das Problem, dass einige Israeliten Leichen-Unrein sind, während das Passah gefeiert werden sollte.
Wie ist mit dieser Situation umzugehen? Können Unreine in dieser speziellen Situation doch am Opfer und Fest teilnehmen?
Ab Vers 9 gibt Jahwe Mose eine positive Antwort: eine Sonderregelung, mit der unter anderem Leichen-Unreine Personen doch am Passah teilhaben können.
In \bibverse{Numbers}{27}{1-11} lesen wir von einem weiteren Sonderfall: Ein Mann stirbt ohne männlichen Erben.
Dürfen seine Töchter stattdessen das Erbe antreten?
Nach diesem Gespräch vor dem Zelt der Begegnung (v2) bringt Mose die Frage vor Gott (v5) und die Antwort ist eine gnädige Sonderfallregelung,
damit die Töchter erben können.

Wenn wir nun zum Römerbrief zurückkehren sehen wir eine sehr ähnliche Funktion, die Jesus Tod erzeugt.
In Christus (durch den Jesus-Christus-Glauben) gibt es eine neue Offenbarung des Gesetzes, die zum Leben führt:
Der wahre Kern, das echte Herz in Gottes Gesetz findet in Jesus Sieg über die Sünde am Kreuz seinen Ausdruck.
Nicht länger ist das Gesetz nur, was Erkenntnis der Sünde bringt (v20), sondern durch die Befähigung in Christus
werden Menschen aus der Sünde erlöst und Tuer des Gesetzes.
Es sind nicht Werke des Gesetzes, sondern der Jesus-Christus-Glauben, der den Menschen befähigt und freisetzt der zu Christus treu ist.
Gleichzeitig rechtfertigt diese Offenbarung auch Gottes eigene Gerechtigkeit:
Wie kann es sein, dass Gott Menschen wie David vergeben hat, wenn doch der Leben hat, der die Torah erfüllt\bibfoot{Leviticus}{18}{5}?
In Christus, dem Ort der neuen Offenbarung, demonstriert Gott, wofür Paulus immer wieder \bibverse{Habakuk}{2}{4} zitiert:
Der durch den Glauben gerechte wird leben!
Diese Botschaft war immer schon wahr,\footnote{
	Das ist kurz gesagt das Argument in Kapitel 4 des Römerbriefes:
	schon bei David und Abraham ist es Glaube, der Sünden vergibt und der Glaube ist es, der dem Menschen als Gerechtigkeit angerechnet wird.
	Wir könnten also sagen: Glaube / Treue = Gerechtigkeit in Gottes Augen; so macht der Glaube Gerecht, weil die beiden Eigenschaften untrennbar zusammen hängen.
	Loyaler Glauben zu Christus ist äquivalent dazu, gerecht zu sein.
	Wer Christus folgt ist automatisch gerecht, und ein Tuer der Torah.
}
aber jetzt durch den Messias gerechtfertigt und allen Menschen offenbart.
Es war nie die Torah, die einen Menschen dazu in die Lage versetzt, gerecht zu sein - es war Glaube (Römer 4).
Und dass diese Tatsache nicht plötzlich Gottes Plan verändert beweist sich darin, wie in Christus jetzt selbst Heiden Tuer der Torah werden (2:14-15).
Das ist das Herz Gottes, die frisch offenbarte Wahrheit in Christus: Gott gibt nicht nur gerechtes Gesetz und setzt einen gerechten Standard (2).
Er selbst greift in seiner selbst aufopfernden Liebe ein, geht über Sünden hinweg, rettet uns aus der Macht der Sünde(3:24b), unter der Juden wie Heiden standen (3:1-20),
und transformiert unser krankes Herz (3:24a).
Der Christus-Glaube, nicht die Werke der Torah, sind es, die einen Menschen gerecht machen.
Und was sagt dieser Christus-Glaube?
Er sagt: alle, die an Gott glauben, gehören zu ihm.
Genau diese Tatsache beweist Gottes Gerechtigkeit in seiner Vergebung für die Sünden in früheren Zeiten.
Schon immer war Gottes Herz Liebe und Vergebung für alle, die ihm folgen - jetzt in Christus endgültig demonstriert und endlich
verbunden mit einer echt wirksamen Transformation des Herzens, die die Werke der Torah nie vollbringen konnten.\footnote{
	Wir denken hier natürlich auch an \bibverse{Ezekiel}{36}{25-27}.
}

\subsection{Römer 5:9-10}
Rom 5:9-10
- Gerecht Gemacht durch Jesus Blut = Transformation = Versöhnung
- Gerettet vor dem Gericht = Gerettet, weil wir tatsächlich durch das Leben in Christus gerecht sind und im Gericht als gerecht befunden werden

\subsection{Römer 5-6 - Der Ungerechtigkeit gestorben, der Gerechtigkeit auferstanden}
tl;dr:
In Römer 5 ab v12 beschreibt Paulus, wie durch Adams Sünde der Tod (=die Verdammnis) in die Welt kam.
Durch Jesus gerechte Tat aber kommt Leben (=Gerecht-Gemacht-Sein; v19).
In Römer 6 beschreibt denselben Kontrast:
Wir sind in Christus gestorben, um damit von der Sünde in der wie gelebt haben gerecht zu sein (v7).
Jetzt könnnen wir in Christus leben. Das passiert, indem wir täglich der Sünde absagen und Werkezeuge der Gerechtigkeit sind, nicht der Ungerechtigkeit.

\subsection{Römer 8:3 - Gottes Urteil über die Sünde}
tl;dr:
Es wird nicht `Jesus für unsere Sünde' verurteilt, sondern `die Sünde im Fleisch'.
Das Ergebnis ist ganz klar, dass die `Rechtsforderung des Gesetzes in uns erfüllt ist', indem wir jetzt durch Loyalität zu Jesus Gerecht handeln -
dadurch, dass wir im Geist leben und die Handlungen des Leibes töten (\bibverse{Romans}{8}{13}).
Es geht nicht darum, wie schuldige Menschen von Gott freigesprochen werden können. Es geht darum, wie die Sünde vernichtet werden kann,
deren Kraft wir Menschen verfallen sind.



\pagebreak
\part{\label{psa-definition}Penal Substitutionary Atonement}
\chapter{Beweistexte}
tl;dr:
Einige Texte klingen sehr stark nach PSA, weil wir PSA gewohnt sind.
Wir gehen die Texte durch und sehen, dass PSA nicht die beste Interpretation ist; gleichzeitig gebe ich jeweils eine alternative Interpretation und zeige, warum sie legitim ist.

\section{Jesaja 52-53 - Partizipation im Leiden des Christus}
tl;dr:
Der leidende Gottesknecht ist kein Opfer.
Er leidet nicht stellvertretend, sondern repräsentativ.
Es gibt kein PSA in \bibverse{Isaiah}{52-53}{}
\subsection{Rezeptionshistorie}
tl;dr:
\bibverse{Isaiah}{52-53}{} wird in der Zeit des zweiten Tempels und in der frühen Kirche konsistent als repräsentatives
Leiden einer Messianischen Figur verstanden.
Es gibt keine frühe Interpretation dieses Textes in einem stellvertretenden, juristischen Rahmen.
Das schließt die Verwendungen dieses Textes im NT ein.

\section{Galater 3:13 - Auslöschung des Bundesfluchs}
- Fluch des Gesetzes ist genau das - Fluch des Gesetzes

TODO - ich bin noch unsicher, wie genau der Bundesfluch zu verstehen ist und wie exakt Jesus damit interagiert

\section{Römer 3:25 - Kein Opfer in Sicht}
tl;dr:
ἱλαστήριον (hilasterion) sollte in diesem Vers nicht als Opfer verstanden werden.
Die Optionen sind (a) ἱλαστήριον = der Deckel der Bundeslade, und zwar \textit{als Ort der Begegnung mit Gott}\Footcite{jesus-as-the-mercy-seat}
(b) ἱλαστήριον als Versöhnungsgeschenk aus der römisch-griechischen Kultur\Footcite{lamb-of-the-free}.

TODO: welche dieser Varianten ich tatsächlich bevorzuge ist noch unklar; ich muss zuerst Bailey vollständig lesen und verstehen.

\section{Römer 8:3 - Gottes Urteil über die Sünde}
tl;dr:
Es wird nicht `Jesus für unsere Sünde' verurteilt, sondern `die Sünde im Fleisch'.

TODO: Die Optionen sind mMn (a) Bait-and-Switch Modell wie in Chrysostomus (b) Jesu Wirken ermöglicht uns, nicht weiter im Fleisch unter der Sünde, sondern im Geist unter Gott zu leben, und das wird als `Verurteilung' bezeichnet.

\section{2. Korinther 5:21}



\chapter{Die Probleme mit Penal Substitutionary Atonement}
\section{Strafverständnis im Tanakh}
tl;dr:
Strafe im Verständnis des Tanakh ist nicht vermittelbar. Es ist immer und ausschließlich die Person die Sündigt, die bestraft werden muss. (Im Gegensatz zu manchen anderen Rechtstexten aus ANE).
Diskussion über Lex Tallionis.
cf. \Cite{old-testament-justice}

\section{Tod ist nicht immer Gottes Strafe}
tl;dr:
Das haben wir schon früher gesehen, aber es ist wichtig genug wiederholt zu werden.

\section{Gott kann wirklich vergeben}
tl;dr:
Gott braucht keine Genugtuung für seinen Zorn oder seine Gerechtigkeit, um zu vergeben.
Gott ist gerecht, wenn er einem bußfertigen Menschen vergibt ohne dafür jemand anderen zu bestrafen.\footnote{
Und so sollen wir auch anderen vergeben.}

\section{Gott ist einer}
tl;dr:
PSA führt zu großen Problem in der Trinität, weil der Sohn vom Vater bestraft wird.

\section{Gott haßt PSA}
tl;dr:
Gott hasst es, den gerechten zu bestrafen oder den ungerechten gerecht zu sprechen (\bibverse{Ezekiel}{18}{}).
Konsistent wird der Sünde für seine Sünde bestraft; es geht nie darum, dass irgendwer für die Sünde bestraft wird.
Sünde kann nicht übertragen werden. Entweder der Sünde trägt die Strafe, oder niemand.

\chapter{Soteriologie in der frühen Kirche}
Die Ansichten zu einer bestimmten Doktrin in der Kirchengeschichte sind nicht unser primärer Weg,
Gottes Sicht auf die Doktrin zu verstehen, aber können uns doch helfen, einige Ansichten als späte Entwicklung auszusortieren.
Es ist mein Standpunkt und meine Argumentation in diesem Kapitel, dass PSA wie in \ref{psa-definition} definiert keine Sicht war, die in den ersten Jahrhunderten der Kirchengeschichte existiert.
Im folgenden werden wir einige Zitate ins Auge nehmen, die \Cite{psa-in-church-history}
verwendet um zu dem Schluss zu kommen ``that the first thousand years of ancient Christianity frequently espoused the teaching that Jesus suffered death, punishment, and a curse for fallen humanity as the penalty for human sin''\Footcite[199, Abstract]{psa-in-church-history}.
Gleichzeitig werde ich - wie in den Kommentaren zu biblischem Material - positiv beschreiben, was diese Väter der Kirche tatsächlich über Jesus' erlösendes Wirken denken.

Bevor wir zu den einzelnen Texten übergehen möchte ich darauf hinweisen, dass PSA für viele seiner Befürworter `der Herz des Evangeliums'\Footcite{psa-is-the-heart-of-the-gospel} ist.
Was erwarten wir in den christlichen Werken der Kirche, wenn das der Fall ist?
Warum lehnen sowohl die Katholische als auch die Orthodoxe Kirche PSA ab, wenn diese Ansicht nicht nur ein Bestandteil, sondern \textit{der} Bestandteil der Christlichen Botschaft ist?
Wenn PSA das Herz des Evangeliums ist, dann erwarten wir in den Schriften der frühen Kirche eine klare Formulierung dieser Doktrin.
Was wir stattdessen finden sollen die nächsten Seiten demonstrieren.

Zunächst findet sich in den Konzilen keine Diskussion und dementsprechend keine Aussagen über die Versöhnungslehre an sich. Gehen wir jetzt also über zu den Aussagen einzelner Kirchenväter
.\footnote{Nicht alle Zitate aus \Cite{psa-in-church-history} sind es überhaupt Wert, diskutiert zu werden.
Als Beispiel nenne ich hier \Cite[49:6]{clement-epistula-ad-corinthios}: ``In Liebe hat der Herr uns angenommen; wegen der Liebe, die er zu uns trug, hat unser Herr Jesus Christus sein Blut hingegeben für uns nach Gottes Willen, sein Fleisch für unser Fleisch, seine Seele für unsere Seelen.''
Nichts in dieser Aussage ist in irgend einer Form distinktiv für PSA.
Christus hat sein Blut und sein Fleisch für uns hingegeben.
Diese Aussage sagt nichts dazu, ob der Vater den Sohn für unsere Sünden bestraft hat.}

\section{Athanasius}
Ich stelle die Diskussion über Athanasius voran, weil sein Verständnis sowohl vom Problem, dass Jesus löst als auch der Lösung, die er dafür gibt leicht verständlich in einem einzelnen Werk dargelegt ist.
Gleichzeitig wird ein gutes Verständnis von Athanasius' Ansicht refokusieren, mit welcher Frage wir tatsächlich an die Texte der Kirchenväter herantreten müssen um zu beantworten, ob sie PSA lehren oder nicht.

In \Cite{de-incarnatione-verbi} gibt Athanasius eine detaillierte Exposition und Verteidigung der Natur der Inkarnation des Wortes, und seines Wirkens.
In diesem Werk können wir also auch erwarten, dass das Erlösungswerk Gottes im Fleischgewordenen Wort zur Sprache kommt und wir werden nicht enttäuscht.
Angefangen von der Erschaffung der Welt erklärt Athanasius, welchen Effekt Sünde hatte, und was das zu lösende Problem seiner Meinung nach ist:

\begin{quote}
Gott hat uns ja nicht nur aus dem Nichts erschaffen, sondern durch die Gnade des Logos
uns auch ein gottgleiches Leben verliehen.
Da aber die Menschen das Ewige von sich wiesen und auf Eingebung des Teufels sich dem Vergänglichen zuwandten, so haben sie selber
damit ihre Vergänglichkeit im Tode verschuldet; sie waren ja wohl, wie vorhin bemerkt,
von Natur aus sterblich, wären aber diesem natürlichen Los dank ihrer Teilnahme am Logos entronnen, wenn sie gut geblieben wären.
Denn um des Logos willen, der mit ihnen
war, hätte die naturgemäße Vernichtung sie nicht betroffen, wie auch die Weisheit sagt:
„Gott hat den Menschen zur Unverweslichkeit erschaffen und als Bild seiner eigenen Ewigkeit;
aber durch den Neid des Teufels kam der Tod in die Welt“\bibfoot{Wisdom}{2}{24}
Wie aber das geschehen war, begann das Sterben der Menschen;
die Vergänglichkeit gewann nunmehr und entfaltete um so stärkere Macht über das ganze Geschlecht denn der natürliche Urstand, als
sie auch die Drohung Gottes wegen der Übertretung des Gebotes ihnen gegenüber voraus hatte.
\attrib{\Cite[5]{de-incarnatione-verbi}}
\end{quote}

Für die Soteriologie in \Cite{de-incarnatione-verbi} ist das zu lösende Problem,
dass sich der Mensch durch seine `Übertretung des Gebotes' von dem natürlichen Zustand
(Unsterblichkeit wegen des geschaffen Seins im Leben des Logos) der Vergänglichkeit (dem Tod) zugewandt hat.

Wir können Athanasius juristische Sprache an folgendem Abschnitt sehr direkt nachvollziehen:
\begin{quote}
So also hat Gott den Menschen erschaffen und ihn in der Unsterblichkeit belassen wollen.
Die Menschen jedoch würdigten den geistigen Verkehr mit Gott wenig, kehrten sich
davon ab, erdachten und ersannen sich die Bosheit, wie im ersten Teil ausgeführt wurde,
und verfielen dem angedrohten Todesurteil. Jetzt sollten sie auch nicht mehr so bleiben,
wie sie geschaffen worden sind, vielmehr sanken sie entsprechend ihrer Denkart immer
tiefer, und der Tod wurde ihr Gewaltherr. Denn die Übertretung des Gebotes warf sie auf
ihren natürlichen Urstand zurück, so daß sie, wie aus dem Nichts geworden, so auch mit
Recht nach Ablauf der Zeit den Verlust ihrer Existenz zu gewärtigen hatten.
Denn wenn
es in ihrer Natur lag, einmal nicht zu sein, und sie erst durch das Eingreifen und die Menschenliebe des Logos ins Dasein gerufen wurden,
so ergab sich als natürliche Folge, daß die Menschen mit dem Verlust ihrer Gottesvorstellung und mit ihrer Abkehr zum Nichtseienden
- nichtseiend ist das Böse, seiend das Gute, weil ja vom Seienden Gott ausgegangen -
auch ihrer ewigen Existenz verlustig gingen, das heißt aber, daß sie der Auflösung anheimfielen und im Tod und in der Verwesung verblieben.
Tatsächlich ist ja der Mensch von
Natur aus sterblich, da er aus dem Nichts entstanden ist. Doch dank seiner Ähnlichkeit
mit dem Seienden hätte er in dem Falle, daß er sie mit einer richtigen Herzensstellung zu
ihm bewahrt hätte, die naturgemäße Auflösung von sich ferngehalten und wäre unverweslich geblieben, wie ja die Weisheit sagt: „Die Beobachtung der Gebote ist die Sicherung der
Unverweslichkeit“\bibfoot{Wisdom}{6}{19}.
\attrib{\Cite[4]{de-incarnatione-verbi}}
\end{quote}

Die Natürliche Ordnung ist, dass der Mensch abseits seiner `richtigen Herzensstellung zu Gott' stirbt.
In Athanasius' Worten ist, das `angedrohte Todesurteil' über die Sünde des Menschen: (a) Unterwerfung unter die Herrschaft des Todes (b) Rückkehr zum natürlichen Zustand abseits des Seienden Logos - Sterblichkeit.
An dieser Stelle wird unmissverständlich klar: für Athanasius ist der Tod ein Feind Gottes, der ihm Diametral gegenübersteht\footnote{nichtseiend ist das Böse, seiend das Gute, weil ja vom Seienden Gott ausgegangen}.
Durch die Erschaffung aus dem Nichts hat das Logos den Menschen Lebendig gemacht.
Durch seine Abwendung von Gott hat der Mensch Tod auf sich gebracht, und zwar \textit{als Folge der natürlichen Ordnung}.
Tod in Athanasius ist keine retributive Strafe Gottes.
Tod ist eine natürlich folgende Konsequenz, die Athanasius im gleichen Atemzug ohne Probleme das `angedrohte Todesurteil' nennen kann.
Wir sehen also, dass das Problem der Tod ist - und zwar Tod als natürliche Ordnung des Menschen abseits der lebensspendenden Einheit mit dem Logos.
Das Problem ist weder Gottes Zorn über Sünde der besänftigt werden muss, noch ein Problem der ausstehenden Strafe Gottes für Sünde, die der Mensch nicht tragen kann.
Das Problem ist, dass der Mensch abseits von Gott Tod ist.
Gleichzeitig ist natürlich genau diese Ordnung von Gott gegeben und der Tod des Menschen, die Übergabe des Menschen an das Nichtssein das aus Sünde resultiert, Gottes Gesetz.
\footnote{So schreibt Athanasius wenig später
``Denn, wie vorhin gesagt, der Tod gewann jetzt über uns gesetzliche Gewalt, und es war unmöglich, dem Gesetze zu entrinnen, weil es von Gott wegen der Sünde erlassen worden war.''
(\Cite[6]{de-incarnatione-verbi}).
Es ist nicht wahr zu sagen, dass Athanasius keine juristische Sprache verwendet, bloß ist nicht Gottes juristische Forderung das Problem, sondern `der Tod gewann die gesetzliche Gewalt':
Der Tod, in Perversion des gerechten Urteil Gottes (Menschen abseits von ihm sterben), kann Sünde verwenden um Tod zu erzeugen.
Die Relationale Trennung des Menschen zu Gott führt zu einer Ontologischen Trennung vom göttlichen Leben des Logos.
Siehe auch \Cite{deification-in-athanasius} für eine deutlich detailliertere Diskussion über die primären Themen in Athanasius' Soteriologie.}

In der Folge macht Athanasius sehr deutlich, warum Buße allein nicht ausreicht, um dieses Problem zu lösen.
(a) Gott wäre nicht wahrhaftig gewesen in der Aufstellung seines Gesetzes, dass Sünde zu Tod führt
(b) Reue hebt den Natürlichen Zustand (also genau dieses gottgegebene, natürliche Gesetz an sich) nicht auf, sondern beendet nur das Sündigen.
Athanasius schreibt ``Wenn also nur Sünde, und nicht auch Vernichtung als deren Folge dagewesen wäre, dann wäre es mit der Reue gut gewesen.''\Footcite[7]{de-incarnatione-verbi}.
In anderen Worten: Wenn nur das Sündigen an sich das Problem wäre, wäre Reue genug um ihr ein Ende zu setzen.
Aber `Vernichtung als deren Folge' (nie in Athanasius als retributive Strafe) kann nicht allein durch Buße beendet werden. Es braucht eine bessere Heilung, die nur das inkarnierte Logos erwirken kann:

\begin{quote}
Er sah das vernünftige Geschlecht zugrunde gehen und den Tod mit seiner Verwesung herrschen über die Menschen;
er sah, wie auch die Strafandrohung für die Sünde uns im Banne des Verderbens festhalte und eine Befreiung daraus vor der Erfüllung des Gesetzes unangebracht wäre;
er sah auch das Unziemliche, das im Falle des Unterganges der Wesen, deren Schöpfer er selber war, gelegen wäre;
er sah auch die überflutende Bosheit der Menschen und wie sie diese nachgerade bis zur Unerträglichkeit zu ihrem eigenen Verderben steigerten;
er sah endlich alle Menschen als die Beute des Todes.
Deshalb erbarmte er sich unseres Geschlechtes, hatte Mitleid mit unserer
Schwachheit, ließ sich herab zu unserer Vergänglichkeit, duldete die Herrschaft des Todes
nicht, und um die Schöpfung gegen den Tod zu schützen und das Werk seines Vaters an
den Menschen nicht vergeblich sein zu lassen, nahm er einen Leib an, und zwar keinen
anderen als den unsrigen.
\attrib{\Cite[8]{de-incarnatione-verbi}}
\end{quote}
Die Probleme sind, in Reihenfolge
(a) die Herrschaft des Todes
(b) die Strafandrohung für Sünde\footnote{
Wie mehrfach gesagt: Die Strafe/das Urteil ist in Athanasius, dass Gott den Menschen seinem natürlichen Nichtssein abseits des lebensspendenden Logos überlässt; die Strafe ist keine retributive Ausschüttung Gottes Zorns über die Sünde.}
- Tod - kann nur aufgehoben werden, wenn das Gesetz erfüllt wird\footnote{Denn ansonsten wäre Gottes Ehre und Wahrhaftigkeit verletzt. Siehe dazu Fußnote 21 in der Zitierten Version.}
(c) Gott will seine eigenen Geschöpfe nicht wieder dem Nichtssein preisgeben
(d) die Bosheit der Menschen führt zu ihrem eigenen Verderben
(e) die Herrschaft des Todes.
Dementsprechend ist auch die Antwort geformt (a) als herabsteigen des Logos zu unserer Vergänglichkeit (b) widerstand gegen und Schutz vor dem Tod (c) damit das Werk des Vaters nicht vergeblich ist.
``Und das tat er aus Liebe zu den Menschen, damit alle in ihm sterben
und so das Gesetz von der Verwesung der Menschen aufgehoben würde, da ja seine Macht
am Leibe des Herrn sich erschöpft hat und bei den gleichartigen Menschen keinen Zugang
mehr finden kann. Auch wollte er die Menschen, die in die Verweslichkeit zurückgefallen
waren, wieder zur Unverweslichkeit erheben und sie vom Tode zu neuem Leben erwecken,
indem er durch die Aneignung des Leibes und die Gnade der Auferstehung den Tod in
ihnen wie eine Stoppel im Feuer vernichtete.''\Footcite[8]{de-incarnatione-verbi}.
Das Gesetz (dass Menschen Sterben) wird aufgehoben, indem sein Tod die `Macht des Gesetzes' erchöpft wird.
Alle sterben in ihm, der Tod kann das Leben des Logos nicht überwinden, und so wird der Tod vernichtet.

\begin{quote}
	Daher hat er den Leib, den er angenommen,
als eine Weihegabe und als ganz makelloses Schlachtopfer in den Tod gegeben und verscheuchte alsobald von allen seinesgleichen den Tod durch
das stellvertretende Opfer.
Denn weil erhaben über alle, hat natürlich der Logos Gottes
mit der Hingabe seines Tempels und der leiblichen Werkstätte ein Entgelt für das Leben
aller entrichtet und die Schuld in seinem Tod bezahlt.
\attrib{\Cite[9]{de-incarnatione-verbi}}
\end{quote}
Somit kommen wir jetzt mit einem klaren Verständnis von Athanasius Problemstellung zu diesen Sätzen, die in \Cite{psa-in-church-history} als Beweis dafür verwendet werden, dass Athanasius PSA lehrt.
Ja, in Athanasius ist Jesus als stellvertretendes Opfer gestorben.
\footnote{τῇ προσφορᾷ τοῦ καταλλήλου NPNF(\Cite{de-incarnatione-verbi-npnf}) übersetzt `by the offering of an equivalent'.
Der Fokus liegt offensichtlicherweise auf der Gleichheit des Opfers, geschaffen durch die Inkarnation.
Die Distinktion zwischen Substitution (Jesus stirbt, damit die Menschheit leben kann) und Repräsentation/Partizipation (alle sterben in ihm, damit das Gesetzt von der Verwesung aufgehoben wird) ist keine, an die Athanasius denkt.
}
Ja, in Athanasius hat Jesus durch seinen Tod die Strafe für die Sünde getragen.
Nein, in keiner Art und Weise lehrt Athanasius PSA.
Wir müssen jeden Author für sich sprechen lassen - in seiner Sprache, über seine Probleme.
Das suchen nach kurzen Zitaten, die unsere Position verteidigen ist außerordentlich gefährlich, führt zu Verwirrung und wir müssen den Denkern der Vergangenheit mit mehr Respekt begegnen.

\Cite{psa-in-church-history} zitiert weiter aus Athanasius' erster Rede gegen die Arianer:
\begin{quote}
	Denn damals wurde die Welt wie ein Schuldiger vom Gesetze gerichtet. Nun aber nahm das Wort das Gericht auf sich, litt im Leibe für alle und schenkte allen das Heil.
	\attrib{\Cite[I.60.]{orationes-contra-arianos}}
\end{quote}
Bevor wir im folgenden betrachten, was Athanasius hier meint, ist direkt offensichtlich, dass dieser Satz nicht verwendet werden kann um PSA zu verteidigen. Es ist in diesem Zitat das Gesetz, dass richtet, nicht der Vater.
Das Gericht des Gesetzes ist, was das Wort auf sich nimmt.
Was also ist das Problem, dass die Inkarnation hier löst?
Im vorhergehenden Absatz lesen wir:
\begin{quote}
Das Gesetz aber wurde von den Engeln verkündet und hat niemand vollkommen gemacht,
da es der Ankunft des Wortes bedurfte, wie Paulus gesagt hat.
Die Ankunft des Wortes aber hat des Vaters Werk vollendet. Und damals herrschte von Adam bis auf Moses der Tod,
das Erscheinen des Wortes aber machte dem Tode ein Ende.
Und wir sterben nicht mehr alle in Adam, sondern werden alle in Christus lebendig gemacht.
Und damals wurde von Dan bis Bersabee das Gesetz verkündigt, und in Judäa allein war Gott bekannt.
Jetzt aber ist über die ganze Erde ihr Ruf gedrungen.
\attrib{\Cite[I.59]{orationes-contra-arianos}}
\end{quote}
In diesem Abschnitt seiner Rede will Athanasius begründen, dass das inkarnierte Wort besser ist als die Engel.
Er vergleicht also das durch die Engel verkündete Gesetzt\bibfoot{Hebrews}{2}{2} mit der Erlösung im inkarnierten Wort.
Was ist es, das in Athanasius' Argument fehlt, und was bewirkt die Inkarnation des Wortes?
Seit Adam herrscht der Tod, und durch das Wort wird dem Tod ein Ende gemacht.
Wieder ist das Problem der Tod, und die Antwort die Inkarnation.
Wir finden auch in dieser Ausführung keinen Bezug dazu, dass Gottes Zorn am Kreuz auf Jesus ausgeschüttet wurde.
Wir finden keine Aussage dazu, dass Jesus für unsere Schuld bestraft wurde.
Auch Vergebung als Konzept spielt absolut keine Rolle.
Was wir stattdessen sehen ist ganz einfach: Das Wort nahm das Gericht (des Gesetzes, das ist der Tod) auf sich, und schenkte allen das Heil (das Leben).
Keine dieser Aussagen ist distinktiv für PSA, und keine der distinktiven Teile von PSA (der Vater bestraft den Sohn für unsere Vergehen) taucht in Athanasius' denken auf.

Wir halten nach dieser längeren Diskussion über Athanasius Soteriologie die folgenden Punkte fest:
(a) wir müssen alle Vorsicht walten lassen, wenn wir die Schriften anderer Kulturen und Zeiten lesen; das selbe Wort kann in einem anderen Kontext eine gänzlich andere Bedeutung haben
(b) die Aussage, die wir tatsächlich in den Kirchenvätern suchen um zu sagen, dass sie PSA verteidigen ist diese: `der Vater hat den Sohn stellvertretend für die Sünde der Menschen bestraft; auf diese Art und Weise bewirkt das Kreuz Vergebung'.
Diese Aussage ist die minimalistische Form von PSA, und diese Aussage werden wir auch im weiteren Verlauf nicht einmal in der Kirche der Antike finden.

\section{Brief an Diognet}
Der Brief \Cite{mathetes-epistula-ad-diognetus} ist laut einigen Verteidigern von PSA `the single best description of penal substitution in the first few centuries, and quite possibly in the history of the church'\Footcite{arnold-psa-in-church-history}, und `if I didn’t know better you could have told me Calvin or Luther wrote this'\Footcite{diognetus-on-christs-atonement}.
Und wenn wir die Stelle in Frage zitieren scheint hier tatsächlich eine wahre Perle der PSA, versteckt im zweiten Nachchristlichen Jahrhundert, zu finden zu sein:
\begin{quote}
Als aber das Mass unserer Ungerechtigkeit voll und es völlig klar geworden war, dass als ihr Lohn Strafe und Tod uns erwarte,
und als der Zeitpunkt gekommen war, den Gott vorausbestimmt hatte, um fortan seine Güte und Macht zu offenbaren,
- o überschwengliche Menschenfreundlichkeit und Liebe Gottes! -
da hasste und verstiess er uns nicht und gedachte nicht des Bösen, sondern war langmütig und geduldig und nahm aus Erbarmen selbst unsere Sünden auf sich;
er selbst gab den eigenen Sohn als Lösepreis für uns, den Heiligen für die Unheiligen,
den Unschuldigen für die Sünder, den Gerechten für die Ungerechten, den Unvergänglichen für die Vergänglichen, den Unsterblichen für die Sterblichen.
Denn was anders war imstande, unsere Sünden zu verdecken als seine Gerechtigkeit?
In wem konnten wir Missetäter und Gottlose gerechtfertigt werden, wenn nicht allein im Sohne Gottes?
Welch süsser Tausch, welch unerforschliches Walten, welch unverhoffte Wohltat,
dass die Ungerechtigkeit vieler in einem Gerechten verborgen würde und die Gerechtigkeit eines einzigen viele Sünder rechtfertige!
\attrib{\Cite[9]{mathetes-epistula-ad-diognetus}}
\end{quote}
Hallelujah! Aber ist das PSA?
In dieser Stelle erscheint das Problem durchaus strafrechtlichen Charakter zu haben (Lohn der Sünde ist Strafe und Tod).
Direkt vor der zitierten Stelle steht weiter, dass wir aus uns selbst nicht in das Reich Gottes gelangen können, aber durch die Kraft Gottes dazu befähigt werden würden.
Hier ist also ganz deutlich ein Problem zu sehen, dass der Mensch aus seiner Natur (seiner Sünde) gerettet werden muss. Wie geschieht diese Erlösung?
`[Gott] nahm aus Erbarmen selbst unsere Sünde auf sich, er selbst gab den eigenen Sohn als Lösepreis für uns', also ist es in diesem Brief \textit{nicht} der Sohn, der 
die Sünden auf sich nahm, denn das selbe Subjekt gibt dann den eigenen Sohn als Lösepreis.
Dieses `die Sünden auf sich nehmen' kann also nicht als Imputation von Schuld verstanden werden, weil dann die Sünde dem Vater imputiert wird.
Stattdessen kann das griechische αὐτὸς τὰς ἡμετέρας ἁμαρτίας ἀνεδέξατο genauso unproblematisch bedeuten
`er nahm die Sünden auf sich: es hat ihn etwas gekostet, sich um sie zu kümmern'\footnote{
	Vergleiche \Cite[ἀναδέχομαι]{lsj}. Die einfache Bedeutung des Wortes `auf sich nehmen' zum Wohle anderer (ohne rechtliche Imputation) ist in dieser Stelle völlig unproblematisch.
	Eine tatsächliches `annehmen der Schuld für ein Vergehen' ist ebenfalls in der semantischen Reichweite dieses Wortes, aber in unserem Fall durch Grammatik und Theologie ausgeschlossen.
} und diese Interpretation ist hier absolut zu bevorzugen, um nicht den Vater als Objekt einer Imputation von Schuld zu verstehen.
Christus ist hier ein Lösepreis für uns - etwas, dass der Vater `hingegeben hat obwohl es etwas kostet', um uns zu befreien.
Er gab `den Heiligen für die Unheiligen', nicht `substitutionell an unserer Stelle', sondern `um uns damit zu nutzen' - `für uns', nicht `an unserer Stelle'.
Auch der tatsächliche Effekt, den der Author hier beschreibt ist nicht, dass die Strafe Gottes für unsere Sünde stattdessen auf Jesus ausgegossen wurde, und so unsere Schuld abgegolten ist.
Der Effekt ist ganz im Gegenteil, unsere Sünden zu verdecken in seiner Gerechtigkeit.

Das Leben und Wirken und Sterben des Sohnes in Gehorsam zum Vater ist die Grundlage, `in dem wir gerechtfertigt werden'.
Die Sprache ist (a) von Gerechtfertigung durch Jesus und (b) der Rettung des Menschen, der aus sich heraus nicht in der Lage ist, zum Leben zu gelangen.
Das Ausgießen der Strafe Gottes für unsere Sünden auf den Sohn ist kein Bestandteil dieser Passage,
sondern die Antwort auf unsere Schuld ist `nichtanrechnen' und `verdecken'.\footnote{
Ich sage nicht, dass PSA mit diesem Text völlig unvereinbar wäre, es ist bloß nicht in dieser Passage zu finden.}

\section{Johannes Chrysostomus}
Mit Johannes Chrysostomus kommen wir jetzt zu der Passage aus den frühen Kirchenvätern, die einer modernen Form von PSA mit Abstand am nächsten kommt.
So lesen wir in seiner Predigt zu \bibverse{IICorinthians}{5}{21}:
\begin{quote}
Denken wir uns, ein König sieht einen Menschen, einen Räuber und Missethäter, den eben die Strafe ereilt,
und der König gibt seinen geliebten, eingebornen, ebenbildlichen Sohn in den Tod und überträgt nebst dem Tode auch die Schuld von jenem Menschen auf den Sohn,
der doch keine Schuld auf sich hat, und befreit so den Verurtheilten nicht bloß von der Strafe,
sondern auch von der Schmach, und erhebt ihn dann noch zu großer Herrschaft: [...]
\attrib{\Cite[11.IV.]{chrysostomus-2-cor}}
\end{quote}
Hier verwendet Chrysostomus eine Analogie, um zu erklären, wie schrecklich es ist nach der Vergebung durch Jesus Wirken weiter zu sündigen.
Seine Analogie ist die, dass der König (Gott der Vater) den Tod und auch die Schuld jedes Menschen auf Jesus legt, um den Menschen zu erhöhen.
Diese Passage klingt tatsächlich stark an PSA an, aber PSA als komplette Theorie ist mit Chrysostomus' gesamtem theologischen Weltbild nicht vereinbar, wie wir an seiner Behandlung von Vers 19 sehen können:
\begin{quote}
	Denn nicht die Thatsache allein ist wunderbar, daß Gott Freund geworden, sondern auch die Weise, wie er es geworden. Und welches ist diese Weise? Indem er ihnen die Sünden erließ; denn anders war es nicht möglich. Darum fährt Paulus fort: „Indem er ihnen nicht anrechnete ihre Übertretungen.“ Hätte uns Gott für die Versündigungen zur Rechenschaft ziehen wollen, so waren wir alle verloren; denn Alle waren gestorben. Aber trotz der Zahl und Größe der Sünden hat Gott statt der Bestrafung sogar sich versöhnt; er hat die Sünden nicht bloß erlassen, sondern gar nicht angerechnet. So müssen denn auch wir den Feinden vergeben, damit wir ebenfalls der gleichen Vergebung theilhaftig werden.
	\attrib{\Cite[11.II]{chrysostomus-2-cor}}
\end{quote}
Gott hat `keine Rechenschaft verlangt', er hat `nicht bestraft', er hat `sie gar nicht angerechnet'.
Alle diese Aussagen stehen im direkten Widerspruch zu PSA, denn dort kann Gott Sünde nicht `nicht anrechnen'.
Die einzige Möglichkeit, damit Gott in PSA Sünden erlassen kann ist, indem er Sünden stattdessen Jesus anrechnet.

In seinem Kommentar zu \bibverse{Gal}{3}{13} lesen wir von Chrysostomus folgendes:
\begin{quote}
Nun war aber doch das Volk jenem anderen Fluche verfallen: „Verflucht sei jeglicher, welcher nicht beharrt bei allem, was geschrieben steht in dem Buche dieses Gesetzes.“
Wie reimt sich das zusammen?
Nun, das Volk war dem Fluche verfallen; denn es beharrte nicht, und es gab keinen, der das ganze Gesetz erfüllt hätte.
Christus aber tauschte sich für diesen Fluch jenen anderen ein, der da sagt: „Verflucht sei jeglicher, welcher am Holze hängt.“
Da nämlich verflucht ist, wer am Holze hängt, und verflucht auch, wer das Gesetz übertritt, so durfte, wer letzteren Fluch lösen wollte, ihm nicht selber verfallen, mußte aber trotzdem einen Fluch auf sich nehmen; deshalb nahm er jenen für diesen auf sich und hat so durch den einen den anderen gelöst.
Und gleichwie für einen zum Tode Verurteilten ein anderer schuldlos den Tod annimmt und ihn der Strafe entreißt, also hat auch Christus getan.
Denn da er nicht dem Fluche (aus) der Übertretung unterworfen war, nahm er statt dessen den Fluch des Kreuzes auf sich, um jenen Fluch zu lösen.
\attrib{\Cite[III.3.]{chrysostomus-gal}}
\end{quote}
Chrysostomus ist hier sehr klar darin, was er unter `er wurde Fluch für uns' versteht.
In seinem denken übernimmt Jesus \textit{nicht} den Fluch des Gesetztes, unter dem der Mensch ist auf eine substitutionelle Weise.
Stattdessen übernimmt er einen \textit{anderen} Fluch, nämlich den des am-Kreuz-hängens\footnote{Die selbe Sicht vertritt \Cite[237f]{paul-and-the-law-in-chrysostom}}.
Auf diese Art und weise stirbt er also zum Wohle des zu Tode Verurteilten, und entreißt ihn von der Strafe (dem Tod): nicht indem er stattdessen dieselbe Strafe auf sich nimmt,
sondern indem er ihn durch das Aufnehmen eines anderen Fluches repräsentiert/austauscht.\footnote{
Ich denke nicht, dass man aus dieser Stelle herauslesen kann, wie diese Repräsentation/Substitution, beziehungsweise das `lösen des Fluches des Gesetzes' für Chrysostomus mechanisch geschieht.
Das Zentrum des Erlösungswerkes in Chrysostomus' Denken ist ganz eindeutig Christus Victor, wie seine Passah-Homilie zeigt.}

Der folgende Text stammt aus einer Homilie zu Himmelfahrt, und hierin finden wir nach meiner Einschätzung den Text in der frühen Kirche, der PSA am nächsten kommt.
Leider ist dieser Text nicht in deutscher Übersetzung vorhanden, weswegen ich ihn hier auf Englisch zitiere.
\begin{quote}
	See then what happens: the mediator is the Son of him who makes consolation, not a human, nor an angel, nor an archangel, nor any servant.
	And what does the mediator do? The work of the mediator.
	Just as when two people turn their backs on each other and do not wish to be reconciled, someone else must come to intervene and break down the enmity between them, so this is what Christ did.
	God was angry with us, we had turned away from God, the Master who loves mankind, and by putting himself in between, Christ reconciled both natures. 

And how did he come between them? He took on himself the punishment that we deserved from the Father and endured the disgrace and insults that we inflicted on God.
Do you desire to learn how he assumed both that punishment from on high and these insults here below? It is said, Christ redeemed us from the curse of the law by becoming a curse for us\bibfoot{Galatians}{3}{13}.
See then how he received the punishment inflicted from above? Consider how he also endured the insults inflicted here below.
The insults of those who insult you, he says, have fallen on me\bibfoot{Psalms}{68}{10}.
See how he dissolved the enmity, how he did not cease to do, suffer, and painstakingly perform all things until he had brought the adversary and enemy back to God himself and made him a friend?
\attrib{\Cite{chrysostomus-fourth-panegyric}}
\end{quote}
Die parallelen zu PSA sind stark, aber nicht exakt: Zunächst ist das Problem tatsächlich Gottes Zorn auf die Menschen, insbesondere sein gerechtes Gericht über Menschliche Sünde.
Jesus tritt explizit in eine Mittlerrolle zwischen einem `angry God' und `turned away humanity'.
Als dieser Mittler nahm er auf sich `the punishment that we deserved from the Father'.
Wir müssen trotzdem vorsichtig damit umgehen, was Chrysostomus hier sagt und was er nicht sagt, denn die Art und Weise, auf die Jesus `the punishment' auf sich nimmt ist explizit so, wie in \bibverse{Gal}{3}{13} geschrieben.
Aus seiner Auslegung dieser selben Stelle wissen wir, wie Chrysostomus die Übernahme des Fluches bzw. der Strafe versteht:
Jesus \textit{starb}, die Strafe die wir gerecht verdient haben. Er war bereit zu sterben, aber nicht indem er \textit{denselben Fluch}, also \textit{die Strafe als Bestrafung} übernommen hat, sondern \textit{einen anderen Fluch}, nämlich \textit{den Tod als äquivalent der Strafe} trug.
In anderen Worten: Chrysostomus sagt nicht, dass Jesus von Gott stellvertretend für uns bestraft wurde.
Er sagt, dass er bereit war, zu sterben - etwas, das wir verdient hätten, nicht er - um Menschen mit Gott zu versöhnen.
Der kleine aber sehr signifikante Unterschied zwischen Chrysostomus' Aussage und PSA ist gerade genau der Kern von PSA:
Hat Jesus diese Versöhnung bewirkt, \textit{indem er von Gott bestraft wurde}?
Chrysostomus' Homilie zu Himmelfahrt \Cite{chrysostomus-fourth-panegyric} beantwortet dieses Frage, wenn wir Chrysostomus Verständnis von \bibverse{Galatians}{3}{13} mit verwenden nicht mit einem Ja.
Um zu demonstrieren, dass mein Verständnis an dieser Stelle keine unnötige Haarspalterei ist sehen wir uns als nächstes Chrysostomus' Kommentar zu \bibverse{Colossians}{2}{13} an:
\begin{quote}
	Also; Wir waren samt und sonders der Sünde und Strafe verfallen; er nahm selbst die Strafe auf sich und hob dadurch Sünde und Strafe auf; die Strafe aber erlitt er am Kreuze.
	Dort nun heftete er den Schuldbrief an; sodann zerriß er ihn, „wie einer, der Macht hat“.
	— Was für einen Schuldbrief? Entweder er meint damit das, was die Israeliten zu Moses sprachen:
	„Alles, was Gott gesagt hat, wollen wir tun und befolgen“; oder wenn nicht dies, daß wir Gott Gehorsam schuldig sind; oder wenn auch das nicht, daß der Teufel den Schuldbrief in Händen hatte, welchen Gott gegen Adam ausstellte, als er sprach:
	„An welchem Tage du von dem Baume issest, wirst du sterben.“
	Diesen Schuldbrief also hatte der Teufel in Händen; und Christus gab ihn uns nicht zurück, sondern riß ihn selbst entzwei zum Zeichen, daß er die Schuld mit Freuden erlasse.
	\attrib{\Cite[6.3]{chrysostomus-col}}
\end{quote}
Auch hier finden wir eine Aussage, die fast wie PSA klingt, aber in Wirklichkeit etwas ganz anderes bedeutet, denn es ist \textit{der Teufel}, der den Schuldschein in der Hand hält.
Das Gesetz ist von Gott gegeben und gut, wie Chrysostomus' immer wieder betont.
Durch unsere Sünden aber führt das Gesetz zu einer Schuld, und in Chrysostomus' Denken ist es \textit{der Teufel, nicht Gott}, der diesen Schuldschein argwöhnisch über unseren Köpfen hält.
Christus nimmt die Strafe (den Tod) auf sich; zerreißt den Schuldschein, und erlässt uns die Schuld mit Freuden.
Nirgendwo finden wir hier die Idee, dass der Vater den Schuldschein über unserem Kopf hält und Christus die gerechte Strafe des Vaters befriedigt, damit der Vater vergeben kann.
Der Teufel ist der Feind, seine Waffe ist die Strafe des Gesetztes, Christus besiegt den Tod und erlässt die Schuld.
Genau darüber beginnt Chrysostomus dann zu reden: ``Während er sich nämlich Hoffnung machte, Christus selbst in seine Gewalt zu bekommen, büßte er sogar alle diejenigen ein, deren er sich schon bemächtigt hatte: als der Leib des Gottmenschen ans Kreuz geschlagen wurde, da standen die Toten auf.''\Footcite[6.3]{chrysostomus-col}

Wieder ein anderes Bild für Jesus Werk am Kreuz finden wir in seiner Homilie zu \bibverse{Romans}{8}{3}:
\begin{quote}
	Solange die Sünde an den Menschen Sünder fand, brachte sie mit Fug und Recht den Tod über sie; als sie aber einmal auf einen Körper traf, der ohne Sünde war, und ihn doch dem Tode überlieferte, da war sie selbst schuldbar geworden, und das Todesurteil wurde über sie gesprochen.
	[...]
	Siehst du da, wie überall die Sünde verurteilt wird, nicht das Fleisch\footnote{Hier: Jesus' menschlicher Körper}? Wie dieses als Sieger erscheint und als Richter über die Sünde
	\attrib{\Cite[14.5.]{chrysostomus-rom}}
\end{quote}
Wieder sehen wir hier das Bild: Christus siegt durch seinen Tod, indem er die Sünde verurteilte (nicht, indem Gott ihn stellvertretend verurteilt).
Die Sünde, der Teufel, der Tod als personifizierter Agent ist das Problem. Christus siegt über sie und bringt Gericht über sie.\footnote{
Diese Idee finden wir auch in anderen Vätern als sogenanntes `bait-and-switch-Modell' (Lockvogelmodell), auch wenn es in verschiedenen Formen ausformuliert ist.}

Zusammenfassend sehen wir, dass der Zorn Gottes über den Menschen in Chrysostomus sehr real ist, und ein echtes Problem.
Jesus tritt als Mediator auf zwischen der sündlosen, heiligen Natur Gottes und der gefallenen, sündigen Natur des Menschen.
Er ist bereit, dafür den Tod (die Strafe des Gesetzes) stellvertretend für Menschen auf sich zu nehmen, \textit{aber nicht indem er statt uns bestraft wurde},
sondern indem er durch den Tod mit dem Menschen eins wird und so die Naturen zusammenführen kann und vor allem:
damit er als Mensch über Sünde und Tod Siegen kann, die die Menschen unterjocht hatten.
Gottes Vergebung aber ist auch in Chrysostomus nicht auf Gottes Todesurteil über Christus als stellvertretendes Opfer prädikiert.

\section{Cyrill}
tl;dr:
Cyrill ist wichtig, weil seine Katechetischen Briefe genau das sind - eine Zusammenfassung des christlichen Glaubens, den neue Konvertiten zu lernen haben.
PSA spielt darin keine Rolle, was ein Problem für die Position darstellt, dass PSA `der Kern des Evangeliums' ist.

% - Cyrill - Catechetical letters

\section{John of Damascus}
tl;dr:
Von Johannes von Damaskus haben wir eine art systematische Theologie, in der PSA wiederum keine Rolle spielt.
% - John of Damascus - The exact exposition of the orthodox faith

\section{Gregory of Nazianzus}
tl;dr:
Gregor von Nazianz hält eine Sicht, die sich als bestes als `Wiederhergestellt Ikone' zusammenfassen lässt.
Jesus heilt unsere korrupte Natur, um uns vollständiger als Imago Dei hinzustellen.

% - Gregory of Naziansus
% 	- Contra Apollinarius Ep.Cl.
% 		- Soteriology = Christology
% 	- Oration 30
% 	- Oration 37
% 	- Oration 39

\section{Eusebeius}
% - Proof of the gospel

\section{Gregory of Nyssa}
% - Gregory of Nyssa

\section{Augustinus}
% 

% kann ich das streichen?
% ich hab keine Lust, cur deus homo zu lesen
\chapter{Penal Substitutionary Atonement in der Kirchengeschichte}
TODO: ggf. fällt diese Kapitel vollständig weg?
\section{Anselm}
tl;dr:
Selbst in Cur Deus Homo finden wir keine exakte Darstellung von PSA, sondern eher eine Satisfaktionstheorie.
% - Cur Deus Homo
\section{Die Reformation}
tl;dr:
Einige Zitate von den Reformatoren die zeigen, wie hier PSA, etwa in seiner modernen Form wirklich auftritt.



\pagebreak
\part{Wenn nicht PSA, was dann?}
\chapter{Ein Kaleidoskop der Versöhnung}
tl;dr:
Jesu Wirken können wir mit einigen verschiedenen Metaphern verstehen, und es hat mehrere Aspekte.
Bloß weil PSA keiner dieser Aspekte ist heißt das nicht, dass Jesu Tod am Kreuz jetzt weniger Gewicht hat.
Hier sind einige der Aspekte, die ich unterstütze.

\section{Christus Victor}
tl;dr:
Jesu Tod und seine Auferstehung siegt über den Tod.

\section{Loskauf}
tl;dr:
Jesu Tod kauft uns frei von der Sünde - wir sind nicht weiter in Sklaverei, sondern jetzt Sklaven des Gottes, der für uns einen Preis bezahlt hat.

\section{Wiederhergestellte Ikone}
tl;dr:
Jesu Blut, ausgeschüttet durch seinen Tod, heilt unsere korrupte Natur.
Wir sind mit Gottes ewigem Leben verbunden indem wir an Jesu Tod und Auferstehung teilhaben.
Diese Verbindung erlaubt uns, Gott wieder vollständig zu repräsentieren.

\section{Apokalyptische Versöhnung}
tl;dr:
Der neue Bund in Christus läutet den Eschaton ein.
Nicht mehr ist der Bund eingeschränkt auf Israeliten - jede Nation hat Zugang ins Allerheiligste durch den neuen Bund in Jesu Blut.

\section{Gegenseitige Partizipation}
tl;dr:
Wir sind nicht aufgerufen uns zurückzulehnen, weil `Jesus gestorben ist, damit wir nicht sterben müssen'.
Wir sind aufgefordert, an Jesu Kreuz und seinem Tod teilzuhaben, da wir wissen, dass wir auch an seiner Auferstehung teilhaben werden.

\chapter{Conclusio}
\section{Warum musste Christus sterben?}
tl;dr:
Jesus musste sterben, weil nur sein Opfer den perfekten neuen Bund einläuten kann.
Jesus musste sterben, damit sein Blut fließt um uns heil zu machen.
Jesus musste sterben, um in allem uns gleich zu sein - nur so kann die göttliche Natur vollständig mit der menschlichen verschmelzen.

\section{Was ist das Evangelium?}
tl;dr:
PSA (und keine Theorie darüber, wie Gott Sünden vergibt) ist kein Teil des Evangeliums.
Das Evangelium ist die frohe Botschaft:
`Jesus der Christus hat seine Königsherrschaft angetreten. Er ist für unsere Sünden gestorben, auferstanden und regiert zur rechten des Vaters.
Wer an ihm teilhat wird mit ihm ewig leben.'

\section{Desiderata - Der Weg nach vorn}
tl;dr:
Ich hoffe, ich konnte zumindest von einigen der folgenden Punkte überzeugen:
\begin{itemize}
	\item PSA ist nicht biblisch, und es ist nicht nötig um die Bibel zu verstehen.
	\item Jesu Tod und Auferstehung erreichen viel mehr, als dass Gott endlich einen Weg gefunden hat, wie er uns vergeben kann.
	\item Christus Victor!
	\item Die frühe Kirche hat nicht zu PSA gehalten.
	\item
		PSA ist nicht der Kern des Evangeliums.
		Unsere Sicht auf die Versöhnungslehre ist eine theologische Debatte, die christliche Gemeinschaft nicht trennen muss.
\end{itemize}

Unser Verständnis der Versöhnungslehre ist nicht zentral, aber trotzdem sehr wichtig und
es hat Auswirkungen darauf, wie wir Gottes Charakter verstehen.
Unsere Partizipation in Christu Wirken ist zentral für praktische Jüngerschaft.

Außerdem ist diese Debatte an vielen Stellen im Protestantismus sehr aufgeheizt und manche Christen können kaum klar darüber denken,
ohne die Gegenseite entweder als Häresie oder als die Anbetung eines kosmischen Kindesmissbrauchs zu verstehen.
Gleichzeitig lässt die Argumentation für PSA
(a) ein Verständnis des Alten Testaments, insbesondere
(b) eine Interaktion mit der modernen Wisschenschaft über die Rituale und Opfer des levitischen Systems, und
(c) eine ehrliche Interakion mit der Kirchengeschichte
vermissen.
Unsere Debattenkultur über theologische Themen muss liebevoller und unsere Argumentation wissenschaftlicher werden.

\appendix
\chapter{\label{app-reparationsopfer}Das Reparationsopfer}
tl;dr:
Das Reinigungsopfer wird für Sakrileg an Sankta dargebracht.
Andere Fälle sind auf diesen zurückzuführen.
Einer dieser Fälle ist der folgende: eine Person enteignet eine andere durch Diebstahl oder ähnliches.
Zu diesem zivilen Vergehen schwört der Täter einen falschen Schwur und belügt somit Gott.
Wenn die Person Reuhe empfindet und die Sünde bekennt, dann kann sie (a) Reparation an die enteignete Person machen und
(b) ein Reparationsopfer bringen, um den Sakrileg des falschen Schwurs auszutilgen.\footnote{
Buße verringert eine bewusste Lüge und falschen Schwur zu einem Vergehen, das wie ein unwissentlich Begangenes entfernt werden kann.
}
Danach vergibt Gott das Vergehen.
Hier ist ein erstaunlicher Fall, wo Buße und eine Reparationszahlung, sowie ein Reparationsopfer als Lösegeld für den Schuldigen Menschen fungieren.
In diesem Fall wird tatäschlich eine Sünde byd rmh durch Buße und ein Opfer vergeben.
Auch in diesem Fall wird keine Schuld auf das Opfer übertragen - es fungiert als Lösegeld, und wird nicht stellvertretend für den Opfernden bestraft.


% Literaturverzeichnis
\clearpage{\pagestyle{empty}\cleardoublepage}
\break
\renewcommand{\thepage}{}
\printbibliography
\printindex

\end{document}
